\chapter{Server Configuration}\label{svn.serverconfig}

A Subversion repository can be accessed simultaneously by clients
running on the same machine on which the repository resides using URLs
carrying the \texttt{file://} scheme. But the typical Subversion setup
involves a single server machine being accessed from clients on
computers all over the office---or, perhaps, all over the world.

This chapter describes how to get your Subversion repository exposed
outside its host machine for use by remote clients. We will cover
Subversion\textquotesingle s currently available server mechanisms,
discussing the configuration and use of each. After reading this
chapter, you should be able to decide which networking setup is right
for your needs, as well as understand how to enable such a setup on your
host computer.

\section{Overview}\label{svn.serverconfig.overview}

Subversion was designed with an abstract repository access layer. This
means that a repository can be programmatically accessed by any sort of
server process, and the client ``repository access'' API allows
programmers to write plug-ins that speak relevant network protocols. In
theory, Subversion can use an infinite number of network
implementations. In practice, there are only two Subversion servers in
widespread use today.

Apache HTTP Server (also known as \texttt{httpd}) is an extremely
popular web server; using the \texttt{mod\_dav\_svn} module, Apache can
access a repository and make it available to clients via the
WebDAV/DeltaV protocol, which is an extension of HTTP. Because Apache is
an extremely extensible server, it provides a number of features ``for
free,'' such as encrypted SSL communication, logging, integration with a
number of third-party authentication systems, and limited built-in web
browsing of repositories.

In the other corner is \texttt{svnserve}: a small, lightweight server
program that speaks a custom protocol with clients. Because its protocol
is explicitly designed for Subversion and is stateful (unlike HTTP), it
provides significantly faster network operations---but at the cost of
some features as well. While it can use SASL to provide a variety of
authentication and encryption options, it has no logging or built-in web
browsing. It is, however, extremely easy to set up and is often the best
option for small teams just starting out with Subversion.

The network protocol which \texttt{svnserve} speaks may also be tunneled
over an SSH connection. This deployment option for \texttt{svnserve}
differs quite a bit in features from a traditional \texttt{svnserve}
deployment. SSH is used to encrypt all communication. SSH is also used
exclusively to authenticate, so real system accounts are required on the
server host (unlike vanilla \texttt{svnserve}, which has its own private
user accounts). Finally, because this setup requires that each user
spawn a private, temporary \texttt{svnserve} process,
it\textquotesingle s equivalent (from a permissions point of view) to
allowing a group of local users to all access the repository via
\texttt{file://} URLs. Path-based access control has no meaning, since
each user is accessing the repository database files directly.

\hyperref[svn.serverconfig.overview.tbl-1]{Comparison of Subversion
server options} provides a quick summary of the three typical server
deployments.

\begin{longtable}[]{@{}>{\raggedright\arraybackslash}p{\dimexpr(\linewidth-2\tabcolsep*4)/4\relax}>{\raggedright\arraybackslash}p{\dimexpr(\linewidth-2\tabcolsep*4)/4\relax}>{\raggedright\arraybackslash}p{\dimexpr(\linewidth-2\tabcolsep*4)/4\relax}>{\raggedright\arraybackslash}p{\dimexpr(\linewidth-2\tabcolsep*4)/4\relax}@{}}
\caption{Comparison of Subversion server
options}\label{svn.serverconfig.overview.tbl-1}\tabularnewline
\toprule\noalign{}
Feature & Apache + mod\_dav\_svn & svnserve & svnserve over SSH \\
\midrule\noalign{}
\endfirsthead
\toprule\noalign{}
Feature & Apache + mod\_dav\_svn & svnserve & svnserve over SSH \\
\midrule\noalign{}
\endhead
\bottomrule\noalign{}
\endlastfoot
Authentication options & HTTP Basic or Digest auth, X.509 certificates,
LDAP, NTLM, or any other mechanism available to Apache httpd & CRAM-MD5
by default; LDAP, NTLM, or any other mechanism available to SASL &
SSH \\
User account options & Private ``users'' file, or other mechanisms
available to Apache httpd (LDAP, SQL, etc.) & Private ``users'' file, or
other mechanisms available to SASL (LDAP, SQL, etc.) & System
accounts \\
Authorization options & Read/write access can be granted over the whole
repository, or specified per path & Read/write access can be granted
over the whole repository, or specified per path & Read/write access
only grantable over the whole repository \\
Encryption & Available via optional SSL (https) & Available via optional
SASL features & Inherent in SSH connection \\
Logging & High-level operational logging of Subversion operations plus
detailed logging at the per-HTTP-request level & High-level operational
logging only & High-level operational logging only \\
Interoperability & Accessible by other WebDAV clients & Talks only to
svn clients & Talks only to svn clients \\
Web viewing & Limited built-in support, or via third-party tools such as
ViewVC & Only via third-party tools such as ViewVC & Only via
third-party tools such as ViewVC \\
Master-slave server replication & Transparent write-proxying available
from slave to master & Can only create read-only slave servers & Can
only create read-only slave servers \\
Speed & Somewhat slower & Somewhat faster & Somewhat faster \\
Initial setup & Somewhat complex & Extremely simple & Moderately
simple \\
\end{longtable}

\section{Choosing a Server
Configuration}\label{svn.serverconfig.choosing}

So, which server should you use? Which is best?

Obviously, there\textquotesingle s no right answer to that question.
Every team has different needs, and the different servers all represent
different sets of trade-offs. The Subversion project itself
doesn\textquotesingle t endorse one server or another, or consider
either server more ``official'' than another.

Here are some reasons why you might choose one deployment over another,
as well as reasons you might \emph{not} choose one.

\subsection{The svnserve
Server}\label{svn.serverconfig.choosing.svnserve}

\begin{description}
\item[Why you might want to use it:]
\hfill
\begin{itemize}
\item
  Quick and easy to set up.
\item
  Network protocol is stateful and noticeably faster than WebDAV.
\item
  No need to create system accounts on server.
\item
  Password is not passed over the network.
\end{itemize}
\item[Why you might want to avoid it:]
\hfill
\begin{itemize}
\item
  By default, only one authentication method is available, the network
  protocol is not encrypted, and the server stores clear text passwords.
  (All these things can be changed by configuring SASL, but
  it\textquotesingle s a bit more work to do.)
\item
  No advanced logging facilities.
\item
  No built-in web browsing. (You\textquotesingle d have to install a
  separate web server and repository browsing software to add this.)
\end{itemize}
\end{description}

\subsection{svnserve over SSH}\label{svn.serverconfig.choosing.svn-ssh}

\begin{description}
\item[Why you might want to use it:]
\hfill
\begin{itemize}
\item
  The network protocol is stateful and noticeably faster than WebDAV.
\item
  You can take advantage of existing SSH accounts and user
  infrastructure.
\item
  All network traffic is encrypted.
\end{itemize}
\item[Why you might want to avoid it:]
\hfill
\begin{itemize}
\item
  Only one choice of authentication method is available.
\item
  No advanced logging facilities.
\item
  It requires users to be in the same system group, or use a shared SSH
  key.
\item
  If used improperly, it can lead to file permission problems.
\end{itemize}
\end{description}

\subsection{The Apache HTTP
Server}\label{svn.serverconfig.choosing.apache}

\begin{description}
\item[Why you might want to use it:]
\hfill
\begin{itemize}
\item
  It allows Subversion to use any of the numerous authentication systems
  already integrated with Apache.
\item
  There is no need to create system accounts on the server.
\item
  Full Apache logging is available.
\item
  Network traffic can be encrypted via SSL.
\item
  HTTP(S) can usually go through corporate firewalls.
\item
  Built-in repository browsing is available via web browser.
\item
  The repository can be mounted as a network drive for transparent
  version control (see
  \hyperref[svn.webdav.autoversioning]{Autoversioning}).
\end{itemize}
\item[Why you might want to avoid it:]
\hfill
\begin{itemize}
\item
  Noticeably slower than \texttt{svnserve}, because HTTP is a stateless
  protocol and requires more network turnarounds.
\item
  Initial setup can be complex.
\end{itemize}
\end{description}

\subsection{Recommendations}\label{svn.serverconfig.choosing.recommendations}

In general, the authors of this book recommend a vanilla
\texttt{svnserve} installation for small teams just trying to get
started with a Subversion server; it\textquotesingle s the simplest to
set up and has the fewest maintenance issues. You can always switch to a
more complex server deployment as your needs change.

Here are some general recommendations and tips, based on years of
supporting users:

\begin{itemize}
\item
  If you\textquotesingle re trying to set up the simplest possible
  server for your group, a vanilla \texttt{svnserve} installation is the
  easiest, fastest route. Note, however, that your repository data will
  be transmitted in the clear over the network. If your deployment is
  entirely within your company\textquotesingle s LAN or VPN, this
  isn\textquotesingle t an issue. If the repository is exposed to the
  wide-open Internet, you might want to make sure that either the
  repository\textquotesingle s contents aren\textquotesingle t sensitive
  (e.g., it contains only open source code), or that you go the extra
  mile in configuring SASL to encrypt network communications.
\item
  If you need to integrate with existing legacy identity systems (LDAP,
  Active Directory, NTLM, X.509, etc.), you must use either the
  Apache-based server or \texttt{svnserve} configured with SASL.
\item
  If you\textquotesingle ve decided to use either Apache or stock
  \texttt{svnserve}, create a single \texttt{svn} user on your system
  and run the server process as that user. Be sure to make the
  repository directory wholly owned by the \texttt{svn} user as well.
  From a security point of view, this keeps the repository data nicely
  siloed and protected by operating system filesystem permissions,
  changeable by only the Subversion server process itself.
\item
  If you have an existing infrastructure that is heavily based on SSH
  accounts, and if your users already have system accounts on your
  server machine, it makes sense to deploy an \texttt{svnserve}-over-SSH
  solution. Otherwise, we don\textquotesingle t widely recommend this
  option to the public. It\textquotesingle s generally considered safer
  to have your users access the repository via (imaginary) accounts
  managed by \texttt{svnserve} or Apache, rather than by full-blown
  system accounts. If your deep desire for encrypted communication still
  draws you to this option, we recommend using Apache with SSL or
  \texttt{svnserve} with SASL encryption instead.
\item
  Do \emph{not} be seduced by the simple idea of having all of your
  users access a repository directly via \texttt{file://} URLs. Even if
  the repository is readily available to everyone via a network share,
  this is a bad idea. It removes any layers of protection between the
  users and the repository: users can accidentally (or intentionally)
  corrupt the repository database, it becomes hard to take the
  repository offline for inspection or upgrade, and it can lead to a
  mess of file permission problems (see
  \hyperref[svn.serverconfig.multimethod]{Supporting Multiple Repository
  Access Methods}). Note that this is also one of the reasons we warn
  against accessing repositories via \texttt{svn+ssh://} URLs---from a
  security standpoint, it\textquotesingle s effectively the same as
  local users accessing via \texttt{file://}, and it can entail all the
  same problems if the administrator isn\textquotesingle t careful.
\end{itemize}

\section{svnserve, a Custom Server}\label{svn.serverconfig.svnserve}

The \texttt{svnserve} program is a lightweight server, capable of
speaking to clients over TCP/IP using a custom, stateful protocol.
Clients contact an \texttt{svnserve} server by using URLs that begin
with the \texttt{svn://} or \texttt{svn+ssh://} scheme. This section
will explain the different ways of running \texttt{svnserve}, how
clients authenticate themselves to the server, and how to configure
appropriate access control to your repositories.

\subsection{Invoking the
Server}\label{svn.serverconfig.svnserve.invoking}

There are a few different ways to run the \texttt{svnserve} program:

\begin{itemize}
\item
  Run \texttt{svnserve} as a standalone daemon, listening for requests.
\item
  Have the Unix \texttt{inetd} daemon temporarily spawn
  \texttt{svnserve} whenever a request comes in on a certain port.
\item
  Have SSH invoke a temporary \texttt{svnserve} over an encrypted
  tunnel.
\item
  Run \texttt{svnserve} as a Microsoft Windows service.
\item
  Run \texttt{svnserve} as a launchd job.
\end{itemize}

The following sections will cover in detail these various deployment
options for \texttt{svnserve}.

\subsubsection{svnserve as
daemon}\label{svn.serverconfig.svnserve.invoking.daemon}

The easiest option is to run \texttt{svnserve} as a standalone
``daemon'' process. Use the \texttt{-d} option for this:

\begin{verbatim}
$ svnserve -d
$               # svnserve is now running, listening on port 3690
\end{verbatim}

When running \texttt{svnserve} in daemon mode, you can use the
\texttt{-\/-listen-port} and \texttt{-\/-listen-host} options to
customize the exact port and hostname to ``bind'' to.

Once we successfully start \texttt{svnserve} as explained previously, it
makes every repository on your system available to the network. A client
needs to specify an \emph{absolute} path in the repository URL. For
example, if a repository is located at \texttt{/var/svn/project1}, a
client would reach it via \url{svn://host.example.com/var/svn/project1}.
To increase security, you can pass the \texttt{-r} option to
\texttt{svnserve}, which restricts it to exporting only repositories
below that path. For example:

\begin{verbatim}
$ svnserve -d -r /var/svn
…
\end{verbatim}

Using the \texttt{-r} option effectively modifies the location that the
program treats as the root of the remote filesystem space. Clients then
use URLs that have that path portion removed from them, leaving much
shorter (and much less revealing) URLs:

\begin{verbatim}
$ svn checkout svn://host.example.com/project1
…
\end{verbatim}

\subsubsection{svnserve via
inetd}\label{svn.serverconfig.svnserve.invoking.inetd}

If you want \texttt{inetd} to launch the process, you need to pass the
\texttt{-i} (\texttt{-\/-inetd}) option. In the following example,
we\textquotesingle ve shown the output from running
\texttt{svnserve\ -i} at the command line, but note that this
isn\textquotesingle t how you actually start the daemon; see the
paragraphs following the example for how to configure \texttt{inetd} to
start \texttt{svnserve}.

\begin{verbatim}
$ svnserve -i
( success ( 2 2 ( ) ( edit-pipeline svndiff1 absent-entries commit-revprops d\
epth log-revprops atomic-revprops partial-replay ) ) )
\end{verbatim}

When invoked with the \texttt{-\/-inetd} option, \texttt{svnserve}
attempts to speak with a Subversion client via \texttt{stdin} and
\texttt{stdout} using a custom protocol. This is the standard behavior
for a program being run via \texttt{inetd}. The IANA has reserved port
3690 for the Subversion protocol, so on a Unix-like system you can add
lines to \texttt{/etc/services} such as these (if they
don\textquotesingle t already exist):

\begin{verbatim}
svn           3690/tcp   # Subversion
svn           3690/udp   # Subversion
\end{verbatim}

If your system is using a classic Unix-like \texttt{inetd} daemon, you
can add this line to \texttt{/etc/inetd.conf}:

\begin{verbatim}
svn stream tcp nowait svnowner /usr/bin/svnserve svnserve -i
\end{verbatim}

Make sure ``svnowner'' is a user that has appropriate permissions to
access your repositories. Now, when a client connection comes into your
server on port 3690, \texttt{inetd} will spawn an \texttt{svnserve}
process to service it. Of course, you may also want to add \texttt{-r}
to the configuration line as well, to restrict which repositories are
exported.

\subsubsection{svnserve via
xinetd}\label{svn.serverconfig.svnserve.invoking.xinetd}

Some operating systems provide the \texttt{xinetd} daemon as an
alternative to \texttt{inetd}. Fortunately, you can configure
\texttt{svnserve} for use with \texttt{xinetd}, too. To do so,
you\textquotesingle ll need to create a configuration file
\texttt{/etc/xinetd.d/svn} with the following contents:

\begin{verbatim}
# default: on
# description: Subversion server for the svn protocol
service svn
{
  disabled        = no
  port            = 3690
  socket_type     = stream
  protocol        = tcp
  wait            = no
  user            = subversion
  server          = /usr/local/bin/svnserve
  server_args     = -i -r /path/to/repositories
}
\end{verbatim}

Be sure that your \texttt{/etc/services} configuration file contains the
definition of the port used for the \texttt{svn} protocol (as described
in \hyperref[svn.serverconfig.svnserve.invoking.inetd]{svnserve via
inetd}), otherwise the daemon will not start correctly.

In Redhat-based distributions, you then need to activate the new service
using \texttt{chkconfig\ -\/-add\ svn}. After doing so, you will be able
to enable and disable the server using the graphical configuration
tools.

\subsubsection{svnserve over a
tunnel}\label{svn.serverconfig.svnserve.invoking.tunnel}

Another way to invoke \texttt{svnserve} is in tunnel mode, using the
\texttt{-t} option. This mode assumes that a remote-service program such
as \texttt{rsh} or \texttt{ssh} has successfully authenticated a user
and is now invoking a private \texttt{svnserve} process \emph{as that
user}. (Note that you, the user, will rarely, if ever, have reason to
invoke \texttt{svnserve} with the \texttt{-t} at the command line;
instead, the SSH daemon does so for you.) The \texttt{svnserve} program
behaves normally (communicating via \texttt{stdin} and \texttt{stdout})
and assumes that the traffic is being automatically redirected over some
sort of tunnel back to the client. When \texttt{svnserve} is invoked by
a tunnel agent like this, be sure that the authenticated user has full
read and write access to the repository database files.
It\textquotesingle s essentially the same as a local user accessing the
repository via \texttt{file://} URLs.

This option is described in much more detail later in this chapter in
\hyperref[svn.serverconfig.svnserve.sshauth]{Tunneling over SSH}.

\subsubsection{svnserve as a Windows
service}\label{svn.serverconfig.svnserve.invoking.winservice}

If your Windows system is a descendant of Windows NT (Windows 2000 or
newer), you can run \texttt{svnserve} as a standard Windows service.
This is typically a much nicer experience than running it as a
standalone daemon with the \texttt{-\/-daemon} (\texttt{-d}) option.
Using daemon mode requires launching a console, typing a command, and
then leaving the console window running indefinitely. A Windows service,
however, runs in the background, can start at boot time automatically,
and can be started and stopped using the same consistent administration
interface as other Windows services.

You\textquotesingle ll need to define the new service using the
command-line tool \texttt{SC.EXE}. Much like the \texttt{inetd}
configuration line, you must specify an exact invocation of
\texttt{svnserve} for Windows to run at startup time:

\begin{verbatim}
C:\> sc create svn
        binpath= "C:\svn\bin\svnserve.exe --service -r C:\repos"
        displayname= "Subversion Server"
        depend= Tcpip
        start= auto
\end{verbatim}

This defines a new Windows service named \texttt{svn} which executes a
particular \texttt{svnserve.exe} command when started (in this case,
rooted at \texttt{C:\textbackslash{}repos}). There are a number of
caveats in the prior example, however.

First, notice that the \texttt{svnserve.exe} program must always be
invoked with the \texttt{-\/-service} option. Any other options to
\texttt{svnserve} must then be specified on the same line, but you
cannot add conflicting options such as \texttt{-\/-daemon}
(\texttt{-d}), \texttt{-\/-tunnel}, or \texttt{-\/-inetd} (\texttt{-i}).
Options such as \texttt{-r} or \texttt{-\/-listen-port} are fine,
though. Second, be careful about spaces when invoking the
\texttt{SC.EXE} command: the \texttt{key=\ value} patterns must have no
spaces between \texttt{key=} and must have exactly one space before the
\texttt{value}. Lastly, be careful about spaces in your command line to
be invoked. If a directory name contains spaces (or other characters
that need escaping), place the entire inner value of \texttt{binpath} in
double quotes, by escaping them:

\begin{verbatim}
C:\> sc create svn
        binpath= "\"C:\program files\svn\bin\svnserve.exe\" --service -r C:\repos"
        displayname= "Subversion Server"
        depend= Tcpip
        start= auto
\end{verbatim}

Also note that the word \texttt{binpath} is misleading---its value is a
\emph{command line}, not the path to an executable.
That\textquotesingle s why you need to surround it with quotes if it
contains embedded spaces.

Once the service is defined, it can be stopped, started, or queried
using standard GUI tools (the Services administrative control panel), or
at the command line:

\begin{verbatim}
C:\> net stop svn
C:\> net start svn
\end{verbatim}

The service can also be uninstalled (i.e., undefined) by deleting its
definition: \texttt{sc\ delete\ svn}. Just be sure to stop the service
first! The \texttt{SC.EXE} program has many other subcommands and
options; run \texttt{sc\ /?} to learn more about it.

\subsubsection{svnserve as a launchd
job}\label{svn.serverconfig.svnserve.invoking.launchd}

Mac OS X (10.4 and higher) uses \texttt{launchd} to manage processes
(including daemons) both system-wide and per-user. A \texttt{launchd}
job is specified by parameters in an XML property list file, and the
\texttt{launchctl} command is used to manage the lifecycle of those
jobs.

When configured to run as a \texttt{launchd} job, \texttt{svnserve} is
automatically launched on demand whenever incoming Subversion
\texttt{svn://} network traffic needs to be handled. This is far more
convenient than a configuration which requires you to manually invoke
\texttt{svnserve} as a long-running background process.

To configure \texttt{svnserve} as a \texttt{launchd} job, first create a
job definition file named
\texttt{/Library/LaunchDaemons/org.apache.subversion.svnserve.plist}.
\hyperref[svn.serverconfig.svnserve.invoking.launchd.ex-1]{A sample
svnserve launchd job definition} provides an example of such a file.

\protect\phantomsection\label{svn.serverconfig.svnserve.invoking.launchd.ex-1}
A sample svnserve launchd job definition

\begin{verbatim}
<?xml version="1.0" encoding="UTF-8"?>
<!DOCTYPE plist PUBLIC "-//Apple//DTD PLIST 1.0//EN"
    "http://www.apple.com/DTDs/PropertyList-1.0.dtd">
<plist version="1.0">
    <dict>
        <key>Label</key>
        <string>org.apache.subversion.svnserve</string>
        <key>ServiceDescription</key>
        <string>Host Subversion repositories using svn:// scheme</string>
        <key>ProgramArguments</key>
        <array>
            <string>/usr/bin/svnserve</string>
            <string>--inetd</string>
            <string>--root=/var/svn</string>
        </array>
        <key>UserName</key>
        <string>svn</string>
        <key>GroupName</key>
        <string>svn</string>
        <key>inetdCompatibility</key>
        <dict>
            <key>Wait</key>
            <false/>
        </dict>
        <key>Sockets</key>
        <dict>
            <key>Listeners</key>
            <array>
                <dict>
                    <key>SockServiceName</key>
                    <string>svn</string>
                    <key>Bonjour</key>
                    <true/>
                </dict>
            </array>
        </dict>
    </dict>
</plist>
\end{verbatim}

The \texttt{launchd} system can be somewhat challenging to learn.
Fortunately, documentation exists for the commands described in this
section. For example, run \texttt{man\ launchd} from the command line to
see the manual page for \texttt{launchd} itself,
\texttt{man\ launchd.plist} to read about the job definition format,
etc.

Once your job definition file is created, you can activate the job using
\texttt{launchctl\ load}:

\begin{verbatim}
$ sudo launchctl load \
       -w /Library/LaunchDaemons/org.apache.subversion.svnserve.plist
\end{verbatim}

To be clear, this action doesn\textquotesingle t actually launch
\texttt{svnserve} yet. It simply tells \texttt{launchd} how to fire up
\texttt{svnserve} when incoming networking traffic arrives on the
\texttt{svn} network port; it will be terminated it after the traffic
has been handled.

Because we want \texttt{svnserve} to be a system-wide daemon process, we
need to use \texttt{sudo} to manage this job as an administrator. Note
also that the \texttt{UserName} and \texttt{GroupName} keys in the
definition file are optional---if omitted, the job will run as the user
who loaded the job.

Deactivating the job is just as easy to do---use
\texttt{launchctl\ unload}:

\begin{verbatim}
$ sudo launchctl unload \
       -w /Library/LaunchDaemons/org.apache.subversion.svnserve.plist
\end{verbatim}

\texttt{launchctl} also provides a way for you to query the status of
jobs. If the job is loaded, there will be line which matches the
\texttt{Label} specified in the job definition file:

\begin{verbatim}
$ sudo launchctl list | grep org.apache.subversion.svnserve
-       0       org.apache.subversion.svnserve
$
\end{verbatim}

\subsection{Built-in Authentication and
Authorization}\label{svn.serverconfig.svnserve.auth}

When a client connects to an \texttt{svnserve} process, the following
things happen:

\begin{itemize}
\item
  The client selects a specific repository.
\item
  The server processes the repository\textquotesingle s
  \texttt{conf/svnserve.conf} file and begins to enforce any
  authentication and authorization policies it describes.
\item
  Depending on the defined policies, one of the following may occur:

  \begin{itemize}
  \item
    The client may be allowed to make requests anonymously, without ever
    receiving an authentication challenge.
  \item
    The client may be challenged for authentication at any time.
  \item
    If operating in tunnel mode, the client will declare itself to be
    already externally authenticated (typically by SSH).
  \end{itemize}
\end{itemize}

The \texttt{svnserve} server, by default, knows only how to send a
CRAM-MD5\footnote{See RFC 2195.} authentication challenge. In essence,
the server sends a small amount of data to the client. The client uses
the MD5 hash algorithm to create a fingerprint of the data and password
combined, and then sends the fingerprint as a response. The server
performs the same computation with the stored password to verify that
the result is identical. \emph{At no point does the actual password
travel over the network.}

If your \texttt{svnserve} server was built with SASL support, it not
only knows how to send CRAM-MD5 challenges, but also likely knows a
whole host of other authentication mechanisms. See
\hyperref[svn.serverconfig.svnserve.sasl]{Using svnserve with SASL}
later in this chapter to learn how to configure SASL authentication and
encryption.

It\textquotesingle s also possible, of course, for the client to be
externally authenticated via a tunnel agent, such as \texttt{ssh}. In
that case, the server simply examines the user it\textquotesingle s
running as, and uses this name as the authenticated username. For more
on this, see the later section,
\hyperref[svn.serverconfig.svnserve.sshauth]{Tunneling over SSH}.

As you\textquotesingle ve already guessed, a
repository\textquotesingle s \texttt{svnserve.conf} file is the central
mechanism for controlling access to the repository. When used in
conjunction with other supplemental files described in this section,
this configuration file offers an administrator a complete solution for
governing user authentication and authorization policies. Each of the
files we\textquotesingle ll discuss uses the format common to other
configuration files (see \hyperref[svn.advanced.confarea]{Runtime
Configuration Area}): section names are marked by square brackets
(\texttt{{[}} and \texttt{{]}}), comments begin with hashes
(\texttt{\#}), and each section contains specific variables that can be
set (\texttt{variable\ =\ value}). Let\textquotesingle s walk through
these files now and learn how to use them.

\subsubsection{Create a users file and
realm}\label{svn.serverconfig.svnserve.auth.users}

For now, the \texttt{{[}general{]}} section of \texttt{svnserve.conf}
has all the variables you need. Begin by changing the values of those
variables: choose a name for a file that will contain your usernames and
passwords and choose an authentication realm:

\begin{verbatim}
[general]
password-db = userfile
realm = example realm
\end{verbatim}

The \texttt{realm} is a name that you define. It tells clients which
sort of ``authentication namespace'' they\textquotesingle re connecting
to; the Subversion client displays it in the authentication prompt and
uses it as a key (along with the server\textquotesingle s hostname and
port) for caching credentials on disk (see
\hyperref[svn.serverconfig.netmodel.credcache]{Caching credentials}).
The \texttt{password-db} variable points to a separate file that
contains a list of usernames and passwords, using the same familiar
format. For example:

\begin{verbatim}
[users]
harry = foopassword
sally = barpassword
\end{verbatim}

The value of \texttt{password-db} can be an absolute or relative path to
the users file. For many admins, it\textquotesingle s easy to keep the
file right in the \texttt{conf/} area of the repository, alongside
\texttt{svnserve.conf}. On the other hand, it\textquotesingle s possible
you may want to have two or more repositories share the same users file;
in that case, the file should probably live in a more public place. The
repositories sharing the users file should also be configured to have
the same realm, since the list of users essentially defines an
authentication realm. Wherever the file lives, be sure to set the
file\textquotesingle s read and write permissions appropriately. If you
know which user(s) \texttt{svnserve} will run as, restrict read access
to the users file as necessary.

\subsubsection{Set access
controls}\label{svn.serverconfig.svnserve.auth.general}

There are two more variables to set in the \texttt{svnserve.conf} file:
they determine what unauthenticated (anonymous) and authenticated users
are allowed to do. The variables \texttt{anon-access} and
\texttt{auth-access} can be set to the value \texttt{none},
\texttt{read}, or \texttt{write}. Setting the value to \texttt{none}
prohibits both reading and writing; \texttt{read} allows read-only
access to the repository, and \texttt{write} allows complete read/write
access to the repository. For example:

\begin{verbatim}
[general]
password-db = userfile
realm = example realm

# anonymous users can only read the repository
anon-access = read

# authenticated users can both read and write
auth-access = write
\end{verbatim}

The example settings are, in fact, the default values of the variables,
should you forget to define them. If you want to be even more
conservative, you can block anonymous access completely:

\begin{verbatim}
[general]
password-db = userfile
realm = example realm

# anonymous users aren't allowed
anon-access = none

# authenticated users can both read and write
auth-access = write
\end{verbatim}

The server process understands not only these ``blanket'' access
controls to the repository, but also finer-grained access restrictions
placed on specific files and directories within the repository. To make
use of this feature, you need to define a file containing more detailed
rules, and then set the \texttt{authz-db} variable to point to it:

\begin{verbatim}
[general]
password-db = userfile
realm = example realm

# Specific access rules for specific locations
authz-db = authzfile
\end{verbatim}

We discuss the syntax of the \texttt{authzfile} file in detail later in
this chapter, in \hyperref[svn.serverconfig.pathbasedauthz]{Path-Based
Authorization}. Note that the \texttt{authz-db} variable
isn\textquotesingle t mutually exclusive with the \texttt{anon-access}
and \texttt{auth-access} variables; if all the variables are defined at
once, \emph{all} of the rules must be satisfied before access is
allowed.

\subsection{Using svnserve with
SASL}\label{svn.serverconfig.svnserve.sasl}

For many teams, the built-in CRAM-MD5 authentication is all they need
from \texttt{svnserve}. However, if your server (and your Subversion
clients) were built with the Cyrus Simple Authentication and Security
Layer (SASL) library, you have a number of authentication and encryption
options available to you.

What Is SASL?

The Cyrus Simple Authentication and Security Layer is open source
software written by Carnegie Mellon University. It adds generic
authentication and encryption capabilities to any network protocol, and
as of Subversion 1.5 and later, both the \texttt{svnserve} server and
\texttt{svn} client know how to make use of this library. It may or may
not be available to you: if you\textquotesingle re building Subversion
yourself, you\textquotesingle ll need to have at least version 2.1 of
SASL installed on your system, and you\textquotesingle ll need to make
sure that it\textquotesingle s detected during
Subversion\textquotesingle s build process. The Subversion command-line
client will report the availability of Cyrus SASL when you run
\texttt{svn\ -\/-version}; if you\textquotesingle re using some other
Subversion client, you might need to check with the package maintainer
as to whether SASL support was compiled in.

SASL comes with a number of pluggable modules that represent different
authentication systems: Kerberos (GSSAPI), NTLM, One-Time-Passwords
(OTP), DIGEST-MD5, LDAP, Secure-Remote-Password (SRP), and others.
Certain mechanisms may or may not be available to you; be sure to check
which modules are provided.

You can download Cyrus SASL (both code and documentation) from
\href{https://www.cyrusimap.org/sasl/}{}.

Normally, when a subversion client connects to \texttt{svnserve}, the
server sends a greeting that advertises a list of the capabilities it
supports, and the client responds with a similar list of capabilities.
If the server is configured to require authentication, it then sends a
challenge that lists the authentication mechanisms available; the client
responds by choosing one of the mechanisms, and then authentication is
carried out in some number of round-trip messages. Even when SASL
capabilities aren\textquotesingle t present, the client and server
inherently know how to use the CRAM-MD5 and ANONYMOUS mechanisms (see
\hyperref[svn.serverconfig.svnserve.auth]{Built-in Authentication and
Authorization}). If server and client were linked against SASL, a number
of other authentication mechanisms may also be available. However,
you\textquotesingle ll need to explicitly configure SASL on the server
side to advertise them.

\subsubsection{Authenticating with
SASL}\label{svn.serverconfig.svnserve.sasl.authn}

To activate specific SASL mechanisms on the server,
you\textquotesingle ll need to do two things. First, create a
\texttt{{[}sasl{]}} section in your repository\textquotesingle s
\texttt{svnserve.conf} file with an initial key-value pair:

\begin{verbatim}
[sasl]
use-sasl = true
\end{verbatim}

Second, create a main SASL configuration file called \texttt{svn.conf}
in a place where the SASL library can find it---typically in the
directory where SASL plug-ins are located. You\textquotesingle ll have
to locate the plug-in directory on your particular system, such as
\texttt{/usr/lib/sasl2/} or \texttt{/etc/sasl2/}. (Note that this is
\emph{not} the \texttt{svnserve.conf} file that lives within a
repository!)

On a Windows server, you\textquotesingle ll also have to edit the system
registry (using a tool such as \texttt{regedit}) to tell SASL where to
find things. Create a registry key named
\texttt{{[}HKEY\_LOCAL\_MACHINE\textbackslash{}SOFTWARE\textbackslash{}Carnegie\ Mellon\textbackslash{}Project\ Cyrus\textbackslash{}SASL\ Library{]}},
and place two keys inside it: a key called \texttt{SearchPath} (whose
value is a path to the directory containing the SASL \texttt{sasl*.dll}
plug-in libraries), and a key called \texttt{ConfFile} (whose value is a
path to the parent directory containing the \texttt{svn.conf} file you
created).

Because SASL provides so many different kinds of authentication
mechanisms, it would be foolish (and far beyond the scope of this book)
to try to describe every possible server-side configuration. Instead, we
recommend that you read the documentation supplied in the \texttt{doc/}
subdirectory of the SASL source code. It goes into great detail about
every mechanism and how to configure the server appropriately for each.
For the purposes of this discussion, we\textquotesingle ll just
demonstrate a simple example of configuring the DIGEST-MD5 mechanism.
For example, if your \texttt{svn.conf} file contains the following:

\begin{verbatim}
pwcheck_method: auxprop
auxprop_plugin: sasldb
sasldb_path: /etc/my_sasldb
mech_list: DIGEST-MD5
\end{verbatim}

you\textquotesingle ve told SASL to advertise the DIGEST-MD5 mechanism
to clients and to check user passwords against a private password
database located at \texttt{/etc/my\_sasldb}. A system administrator can
then use the \texttt{saslpasswd2} program to add or modify usernames and
passwords in the database:

\begin{verbatim}
$ saslpasswd2 -c -f /etc/my_sasldb -u realm username
\end{verbatim}

A few words of warning: first, make sure the ``realm'' argument to
\texttt{saslpasswd2} matches the same realm you\textquotesingle ve
defined in your repository\textquotesingle s \texttt{svnserve.conf}
file; if they don\textquotesingle t match, authentication will fail.
Also, due to a shortcoming in SASL, the common realm must be a string
with no space characters. Finally, if you decide to go with the standard
SASL password database, make sure the \texttt{svnserve} program has read
access to the file (and possibly write access as well, if
you\textquotesingle re using a mechanism such as OTP).

This is just one simple way of configuring SASL. Many other
authentication mechanisms are available, and passwords can be stored in
other places such as in LDAP or a SQL database. Consult the full SASL
documentation for details.

Remember that if you configure your server to only allow certain SASL
authentication mechanisms, this forces all connecting clients to have
SASL support as well. Any Subversion client built without SASL support
(which includes all pre-1.5 clients) will be unable to authenticate. On
the one hand, this sort of restriction may be exactly what you want
(``My clients must all use Kerberos!''). However, if you still want
non-SASL clients to be able to authenticate, be sure to advertise the
CRAM-MD5 mechanism as an option. All clients are able to use CRAM-MD5,
whether they have SASL capabilities or not.

\subsubsection{SASL
encryption}\label{svn.serverconfig.svnserve.sasl.encryption}

SASL is also able to perform data encryption if a particular mechanism
supports it. The built-in CRAM-MD5 mechanism doesn\textquotesingle t
support encryption, but DIGEST-MD5 does, and mechanisms such as SRP
actually require use of the OpenSSL library. To enable or disable
different levels of encryption, you can set two values in your
repository\textquotesingle s \texttt{svnserve.conf} file:

\begin{verbatim}
[sasl]
use-sasl = true
min-encryption = 128
max-encryption = 256
\end{verbatim}

The \texttt{min-encryption} and \texttt{max-encryption} variables
control the level of encryption demanded by the server. To disable
encryption completely, set both values to 0. To enable simple
checksumming of data (i.e., prevent tampering and guarantee data
integrity without encryption), set both values to 1. If you wish to
allow---but not require---encryption, set the minimum value to 0, and
the maximum value to some bit length. To require encryption
unconditionally, set both values to numbers greater than 1. In our
previous example, we require clients to do at least 128-bit encryption,
but no more than 256-bit encryption.

\subsection{Tunneling over SSH}\label{svn.serverconfig.svnserve.sshauth}

\texttt{svnserve}\textquotesingle s built-in authentication (and SASL
support) can be very handy, because it avoids the need to create real
system accounts. On the other hand, some administrators already have
well-established SSH authentication frameworks in place. In these
situations, all of the project\textquotesingle s users already have
system accounts and the ability to ``SSH into'' the server machine.

It\textquotesingle s easy to use SSH in conjunction with
\texttt{svnserve}. The client simply uses the \texttt{svn+ssh://} URL
scheme to connect:

\begin{verbatim}
$ whoami
harry

$ svn list svn+ssh://host.example.com/repos/project
harryssh@host.example.com's password:  *****

foo
bar
baz
…
\end{verbatim}

In this example, the Subversion client is invoking a local \texttt{ssh}
process, connecting to \texttt{host.example.com}, authenticating as the
user \texttt{harryssh} (according to SSH user configuration), then
spawning a private \texttt{svnserve} process on the remote machine
running as the user \texttt{harryssh}. The \texttt{svnserve} command is
being invoked in tunnel mode (\texttt{-t}), and its network protocol is
being ``tunneled'' over the encrypted connection by \texttt{ssh}, the
tunnel agent. If the client performs a commit, the authenticated
username \texttt{harryssh} will be used as the author of the new
revision.

The important thing to understand here is that the Subversion client is
\emph{not} connecting to a running \texttt{svnserve} daemon. This method
of access doesn\textquotesingle t require a daemon, nor does it notice
one if present. It relies wholly on the ability of \texttt{ssh} to spawn
a temporary \texttt{svnserve} process, which then terminates when the
network connection is closed.

When using \texttt{svn+ssh://} URLs to access a repository, remember
that it\textquotesingle s the \texttt{ssh} program prompting for
authentication, and \emph{not} the \texttt{svn} client program. That
means there\textquotesingle s no automatic password-caching going on
(see \hyperref[svn.serverconfig.netmodel.credcache]{Caching
credentials}). The Subversion client often makes multiple connections to
the repository, though users don\textquotesingle t normally notice this
due to the password caching feature. When using \texttt{svn+ssh://}
URLs, however, users may be annoyed by \texttt{ssh} repeatedly asking
for a password for every outbound connection. The solution is to use a
separate SSH password-caching tool such as \texttt{ssh-agent} on a
Unix-like system, or \texttt{pageant} on Windows.

When running over a tunnel, authorization is primarily controlled by
operating system permissions to the repository\textquotesingle s
database files; it\textquotesingle s very much the same as if Harry were
accessing the repository directly via a \texttt{file://} URL. If
multiple system users are going to be accessing the repository directly,
you may want to place them into a common group, and
you\textquotesingle ll need to be careful about umasks (be sure to read
\hyperref[svn.serverconfig.multimethod]{Supporting Multiple Repository
Access Methods} later in this chapter). But even in the case of
tunneling, you can still use the \texttt{svnserve.conf} file to block
access, by simply setting \texttt{auth-access\ =\ read} or
\texttt{auth-access\ =\ none}.\footnote{Note that using any sort of
  \texttt{svnserve}-enforced access control at all only makes sense if
  the users cannot bypass it and access the repository directory
  directly using other tools (such as \texttt{cd} and \texttt{vi});
  implementing such restrictions is described in
  \hyperref[svn.serverconfig.svnserve.sshtricks.fixedcmd]{Controlling
  the invoked command}.}

You\textquotesingle d think that the story of SSH tunneling would end
here, but it doesn\textquotesingle t. Subversion allows you to create
custom tunnel behaviors in your runtime \texttt{config} file (see
\hyperref[svn.advanced.confarea]{Runtime Configuration Area}). For
example, suppose you want to use RSH instead of SSH.\footnote{We
  don\textquotesingle t actually recommend this, since RSH is notably
  less secure than SSH.} In the \texttt{{[}tunnels{]}} section of your
\texttt{config} file, simply define it like this:

\begin{verbatim}
[tunnels]
rsh = rsh --
\end{verbatim}

And now, you can use this new tunnel definition by using a URL scheme
that matches the name of your new variable:
\texttt{svn+rsh://host/path}. When using the new URL scheme, the
Subversion client will actually be running the command
\texttt{rsh\ -\/-\ host\ svnserve\ -t} behind the scenes. If you include
a username in the URL (e.g., \texttt{svn+rsh://username@host/path}), the
client will also include that in its command
(\texttt{rsh\ -\/-\ username@host\ svnserve\ -t}).

Notice that when defining an RSH-based tunnel, we\textquotesingle ve
added the \texttt{-\/-} end-of-options argument to the tunnel command
line. This is to prevent a malformed hostname from being treated as
another option to the tunnel command. You should do the same for other
tunnel programs (for example, SSH).

But you can define new tunneling schemes to be much more clever than
that:

\begin{verbatim}
[tunnels]
joessh = $JOESSH /opt/alternate/ssh -p 29934 --
\end{verbatim}

This example demonstrates a couple of things. First, it shows how to
make the Subversion client launch a very specific tunneling binary (the
one located at \texttt{/opt/alternate/ssh}) with specific options. In
this case, accessing an \texttt{svn+joessh://} URL would invoke the
particular SSH binary with \texttt{-p\ 29934} as arguments---useful if
you want the tunnel program to connect to a nonstandard port.

Second, it shows how to define a custom environment variable that can
override the name of the tunneling program. Setting the
\texttt{SVN\_SSH} environment variable is a convenient way to override
the default SSH tunnel agent. But if you need to have several different
overrides for different servers, each perhaps contacting a different
port or passing a different set of options to SSH, you can use the
mechanism demonstrated in this example. Now if we were to set the
\texttt{JOESSH} environment variable, its value would override the
entire value of the tunnel variable---\texttt{\$JOESSH} would be
executed instead of \texttt{/opt/alternate/ssh\ -p\ 29934}.

\subsection{SSH Configuration
Tricks}\label{svn.serverconfig.svnserve.sshtricks}

It\textquotesingle s possible to control not only the way in which the
client invokes \texttt{ssh}, but also to control the behavior of
\texttt{sshd} on your server machine. In this section,
we\textquotesingle ll show how to control the exact \texttt{svnserve}
command executed by \texttt{sshd}, as well as how to have multiple users
share a single system account.

\subsubsection{Initial
setup}\label{svn.serverconfig.svnserve.sshtricks.setup}

To begin, locate the home directory of the account
you\textquotesingle ll be using to launch \texttt{svnserve}. Make sure
the account has an SSH public/private keypair installed, and that the
user can log in via public-key authentication. Password authentication
will not work, since all of the following SSH tricks revolve around
using the SSH \texttt{authorized\_keys} file.

If it doesn\textquotesingle t already exist, create the
\texttt{authorized\_keys} file (on Unix, typically
\texttt{\textasciitilde{}/.ssh/authorized\_keys}). Each line in this
file describes a public key that is allowed to connect. The lines are
typically of the form:

\begin{verbatim}
  ssh-dsa AAAABtce9euch… user@example.com
\end{verbatim}

The first field describes the type of key, the second field is the
base64-encoded key itself, and the third field is a comment. However,
it\textquotesingle s a lesser known fact that the entire line can be
preceded by a \texttt{command} field:

\begin{verbatim}
  command="program" ssh-dsa AAAABtce9euch… user@example.com
\end{verbatim}

When the \texttt{command} field is set, the SSH daemon will run the
named program instead of the typical tunnel-mode \texttt{svnserve}
invocation that the Subversion client asks for. This opens the door to a
number of server-side tricks. In the following examples, we abbreviate
the lines of the file as:

\begin{verbatim}
  command="program" TYPE KEY COMMENT
\end{verbatim}

\subsubsection{Controlling the invoked
command}\label{svn.serverconfig.svnserve.sshtricks.fixedcmd}

Because we can specify the executed server-side command,
it\textquotesingle s easy to name a specific \texttt{svnserve} binary to
run and to pass it extra arguments:

\begin{verbatim}
  command="/path/to/svnserve -t -r /virtual/root" TYPE KEY COMMENT
\end{verbatim}

In this example, \texttt{/path/to/svnserve} might be a custom wrapper
script around \texttt{svnserve} which sets the umask (see
\hyperref[svn.serverconfig.multimethod]{Supporting Multiple Repository
Access Methods}). It also shows how to anchor \texttt{svnserve} in a
virtual root directory, just as one often does when running
\texttt{svnserve} as a daemon process. This might be done either to
restrict access to parts of the system, or simply to relieve the user of
having to type an absolute path in the \texttt{svn+ssh://} URL.

It\textquotesingle s also possible to have multiple users share a single
account. Instead of creating a separate system account for each user,
generate a public/private key pair for each person. Then place each
public key into the \texttt{authorized\_keys} file, one per line, and
use the \texttt{-\/-tunnel-user} option:

\begin{verbatim}
  command="svnserve -t --tunnel-user=harry" TYPE1 KEY1 harry@example.com
  command="svnserve -t --tunnel-user=sally" TYPE2 KEY2 sally@example.com
\end{verbatim}

This example allows both Harry and Sally to connect to the same account
via public key authentication. Each of them has a custom command that
will be executed; the \texttt{-\/-tunnel-user} option tells
\texttt{svnserve} to assume that the named argument is the authenticated
user. Without \texttt{-\/-tunnel-user}, it would appear as though all
commits were coming from the one shared system account.

A final word of caution: giving a user access to the server via
public-key in a shared account might still allow other forms of SSH
access, even if you\textquotesingle ve set the \texttt{command} value in
\texttt{authorized\_keys}. For example, the user may still get shell
access through SSH or be able to perform X11 or general port forwarding
through your server. To give the user as little permission as possible,
you may want to specify a number of restrictive options immediately
after the \texttt{command}:

\begin{verbatim}
  command="svnserve -t --tunnel-user=harry",no-port-forwarding,no-agent-forw
arding,no-X11-forwarding,no-pty TYPE1 KEY1 harry@example.com
\end{verbatim}

Note that this all must be on one line---truly on one line---since SSH
\texttt{authorized\_keys} files do not even allow the conventional
backslash character (\texttt{\textbackslash{}}) for line continuation.
The only reason we\textquotesingle ve shown it with a line break is to
fit it on the physical page of a book.

\subsection{svnserve Configuration
Reference}\label{svn.serverconfig.svnserve.ref}

In the previous sections, we\textquotesingle ve mentioned numerous
configuration options that administrators can use in their
\texttt{svnserve.conf} files to configure the behavior of Subversion as
accessed via Subversion\textquotesingle s \texttt{svnserve} server
option. In this section, we\textquotesingle ll quickly summarize
\emph{all} the configuration options supported by this server.

The \texttt{svnserve.conf} configuration file uses a typical INI-style
format, with name/value pairs of options grouped into named sections.
(This is conveniently the same format used by
Subversion\textquotesingle s runtime configuration area on the client
side of the network.) We\textquotesingle ll describe herein each of
those named sections and the options available for use within them.

By default, \texttt{svnserve} will consult per-repository configuration
files located at \texttt{conf/svnserve.conf} within the physical
directory structure of the repository. To instead use a single
configuration file whose values apply to all repositories served via an
instance of \texttt{svnserve}, use the \texttt{-\/-config-file} option
when starting your server.

In the following sections, we will refer to the \texttt{svnserve}
configuration file by its canonical name, \texttt{svnserve.conf}. The
filename of actual configuration file used by your \texttt{svnserve}
instance might be something else, though. We trust this
won\textquotesingle t be too confusing.

\subsubsection{General
configuration}\label{svn.serverconfig.svnserve.ref.general}

The \texttt{{[}general{]}} section contains the most commonly used and
broadly focused \texttt{svnserve} configuration options.

\begin{description}
\item[\texttt{anon-access}]
Controls the access level granted to unauthenticated (anonymous) users.
Valid values are \texttt{write}, \texttt{read}, and \texttt{none}, with
\texttt{read} being the default value.
\item[\texttt{auth-access}]
Controls the access level granted to authenticaed users. Valid values
are \texttt{write}, \texttt{read}, and \texttt{none}, with
\texttt{write} being the default value.
\item[\texttt{authz-db}]
Specifies the location of the repository access file as described in
\hyperref[svn.serverconfig.pathbasedauthz.getting-started]{Getting
Started with Path-Based Access Control}. If a regular local path is
used, then unless that path begins with a forward-slash character
(\texttt{/}), it is interpreted as a path relative to the directory
containing the \texttt{svnserve.conf} configuration file. If no path is
specified, path-based access control will be disabled.

As a special consideration, you may also specify the location of an
access file which is versioned inside a Subversion repository. Use a
local URL (one which begins with \texttt{file://}) to refer to an
absolute Subversion-versioned access file. Alternatively, use a
repository relative URL (one which begins with \texttt{\^{}/}) to cause
\texttt{svnserve} to consult for each repository the access file stored
at the specified relative URL within that repository.
\item[\texttt{force-username-case}]
Specifies the case normalization applied to usernames before comparing
them against the rules found in the access file (specified by the
\texttt{authz-db} option). Valid values are \texttt{upper} (to uppercase
the usernames), \texttt{lower} (to lowercase the usernames), and
\texttt{none} (to perform no normalization at all). By default,
\texttt{svnserve} will not perform any case normalization on usernames.
\item[\texttt{groups-db}]
Specifies the path of the groups file. If a regular local path is used,
then unless that path begins with a forward-slash character
(\texttt{/}), it is interpreted as a path relative to the directory
containing the \texttt{svnserve.conf} configuration file.

You may also specify the location of a groups file which is versioned
inside a Subversion repository. Use a local URL (one which begins with
\texttt{file://}) to refer to an absolute Subversion-versioned file.
Alternatively, use a repository relative URL (one which begins with
\texttt{\^{}/}) to cause \texttt{svnserve} to consult for each
repository the group file stored at the specified relative URL within
that repository.
\item[\texttt{hooks-env}]
Specifies the path to the hook script environment configuration file.
This option overrides the per-repository default location for this file,
and can be used to configure the hook script environment for multiple
repositories in a single file if an absolute path is specified. Unless
you specify an absolute path, the file\textquotesingle s location is
interpreted as relative to the directory containing the
\texttt{svnserve.conf} configuration file.

See \hyperref[svn.reposadmin.hooks.configuration]{Hook script
environment configuration} for detailed information regarding the hook
script environment configuration file.
\item[\texttt{password-db}]
Specifies the path of the password database file. Unless the path
specified begins with a forward-slash character (\texttt{/}), it is
interpreted as a path relative to the directory containing the
\texttt{svnserve.conf} configuration file. Note that if the SASL feature
is used, this option will be ignored.
\item[\texttt{realm}]
Specifies the authentication realm of the repository. This is primarily
used by the client to associate cached authentication credentials with a
specific repository or set of repositories. As such, it is best that the
specified realm be unique across your repositories unless those
repositories share the same password database. By default, the
repository\textquotesingle s UUID is used as its authentication realm.
\end{description}

\subsubsection{Cyrus SASL
configuration}\label{svn.serverconfig.svnserve.ref.sasl}

The \texttt{{[}sasl{]}} section contains configuration which is specific
to the optional Cyrus Simple Authentication and Security Layer (SASL)
integration feature of \texttt{svnserve}. See
\hyperref[svn.serverconfig.svnserve.sasl]{Using svnserve with SASL} for
a more thorough description of this feature and the benefits it
provides.

\begin{description}
\item[\texttt{max-encryption}]
Specifies---as an integer bit-width---the maximum desired strength of
the security layer\textquotesingle s encryption algorithm. The special
value \texttt{0} means "no encryption", and the special value \texttt{1}
means "integrity checking only". The default value for this option is
\texttt{256} (256-bit encryption).
\item[\texttt{min-encryption}]
Specifies---as an integer bit-width---the minimum desired strength of
the security layer\textquotesingle s encryption algorithm. The special
value \texttt{0} means "no encryption", and the special value \texttt{1}
means "integrity checking only". The default value for this option is
\texttt{0} (no encryption).
\item[\texttt{use-sasl}]
Specifies (as a \texttt{true} or \texttt{false} value) whether to enable
the Cyrus SASL feature. Note that this feature is only available if
\texttt{svnserve} was built with support for the feature. This feature
is disabled by default.
\end{description}

\section{httpd, the Apache HTTP Server}\label{svn.serverconfig.httpd}

The Apache HTTP Server is a ``heavy-duty'' network server that
Subversion can leverage. Via a custom module, \texttt{httpd} makes
Subversion repositories available to clients via the
WebDAV/DeltaV\footnote{See \href{http://www.webdav.org/}{}.} protocol,
which is an extension to HTTP 1.1. This protocol takes the ubiquitous
HTTP protocol that is the core of the World Wide Web, and adds
writing---specifically, versioned writing---capabilities. The result is
a standardized, robust system that is conveniently packaged as part of
the Apache 2.0 software, supported by numerous operating systems and
third-party products, and doesn\textquotesingle t require network
administrators to open up yet another custom port.\footnote{They really
  hate doing that.} While an Apache-Subversion server has more features
than \texttt{svnserve}, it\textquotesingle s also a bit more difficult
to set up. With flexibility often comes more complexity.

Much of the following discussion includes references to Apache
configuration directives. While some examples are given of the use of
these directives, describing them in full is outside the scope of this
chapter. The Apache team maintains excellent documentation, publicly
available on their web site at \href{https://httpd.apache.org}{}. For
example, a general reference for the configuration directives is located
at \href{https://httpd.apache.org/docs/current/mod/directives.html}{}.

Also, as you make changes to your Apache setup, it is likely that
somewhere along the way a mistake will be made. If you are not already
familiar with Apache\textquotesingle s logging subsystem, you should
become aware of it. In your \texttt{httpd.conf} file are directives that
specify the on-disk locations of the access and error logs generated by
Apache (the \texttt{CustomLog} and \texttt{ErrorLog} directives,
respectively). Subversion\textquotesingle s \texttt{mod\_dav\_svn} uses
Apache\textquotesingle s error logging interface as well. You can always
browse the contents of those files for information that might reveal the
source of a problem that is not clearly noticeable otherwise.

\subsection{Prerequisites}\label{svn.serverconfig.httpd.prereqs}

To network your repository over HTTP, you basically need four
components, available in two packages. You\textquotesingle ll need
Apache \texttt{httpd} 2.0 or newer, the \texttt{mod\_dav} DAV module
that comes with it, Subversion, and the \texttt{mod\_dav\_svn}
filesystem provider module distributed with Subversion. Once you have
all of those components, the process of networking your repository is as
simple as:

\begin{itemize}
\item
  Getting httpd up and running with the \texttt{mod\_dav} module
\item
  Installing the \texttt{mod\_dav\_svn} backend to \texttt{mod\_dav},
  which uses Subversion\textquotesingle s libraries to access the
  repository
\item
  Configuring your \texttt{httpd.conf} file to export (or expose) the
  repository
\end{itemize}

You can accomplish the first two items either by compiling
\texttt{httpd} and Subversion from source code or by installing prebuilt
binary packages of them on your system. For the most up-to-date
information on how to compile Subversion for use with the Apache HTTP
Server, as well as how to compile and configure Apache itself for this
purpose, see the \texttt{INSTALL} file in the top level of the
Subversion source code tree.

\subsection{Basic Apache
Configuration}\label{svn.serverconfig.httpd.basic}

Once you have all the necessary components installed on your system, all
that remains is the configuration of Apache via its \texttt{httpd.conf}
file. Instruct Apache to load the \texttt{mod\_dav\_svn} module using
the \texttt{LoadModule} directive. This directive must precede any other
Subversion-related configuration items. If your Apache was installed
using the default layout, your \texttt{mod\_dav\_svn} module should have
been installed in the \texttt{modules} subdirectory of the Apache
install location (often \texttt{/usr/local/apache2}). The
\texttt{LoadModule} directive has a simple syntax, mapping a named
module to the location of a shared library on disk:

\begin{verbatim}
LoadModule dav_svn_module     modules/mod_dav_svn.so
\end{verbatim}

Apache interprets the \texttt{LoadModule} configuration
item\textquotesingle s library path as relative to its own server root.
If configured as previously shown, Apache will look for the Subversion
DAV module shared library in its own \texttt{modules/} subdirectory.
Depending on how Subversion was installed on your system, you might need
to specify a different path for this library altogether, perhaps even an
absolute path such as in the following example:

\begin{verbatim}
LoadModule dav_svn_module     C:/Subversion/libexec/mod_dav_svn.so
\end{verbatim}

Note that if \texttt{mod\_dav} was compiled as a shared object (instead
of statically linked directly to the \texttt{httpd} binary),
you\textquotesingle ll need a similar \texttt{LoadModule} statement for
it, too. Be sure that it comes before the \texttt{mod\_dav\_svn} line:

\begin{verbatim}
LoadModule dav_module         modules/mod_dav.so
LoadModule dav_svn_module     modules/mod_dav_svn.so
\end{verbatim}

At a later location in your configuration file, you now need to tell
Apache where you keep your Subversion repository (or repositories). The
\texttt{Location} directive has an XML-like notation, starting with an
opening tag and ending with a closing tag, with various other
configuration directives in the middle. The purpose of the
\texttt{Location} directive is to instruct Apache to do something
special when handling requests that are directed at a given URL or one
of its children. In the case of Subversion, you want Apache to simply
hand off support for URLs that point at versioned resources to the DAV
layer. You can instruct Apache to delegate the handling of all URLs
whose path portions (the part of the URL that follows the
server\textquotesingle s name and the optional port number) begin with
\texttt{/repos/} to a DAV provider whose repository is located at
\texttt{/var/svn/repository} using the following \texttt{httpd.conf}
syntax:

\begin{verbatim}
<Location /repos>
  DAV svn
  SVNPath /var/svn/repository
</Location>
\end{verbatim}

If you plan to support multiple Subversion repositories that will reside
in the same parent directory on your local disk, you can use an
alternative directive---\texttt{SVNParentPath}---to indicate that common
parent directory. For example, if you know you will be creating multiple
Subversion repositories in a directory \texttt{/var/svn} that would be
accessed via URLs such as \url{http://my.server.com/svn/repos1},
\url{http://my.server.com/svn/repos2}, and so on, you could use the
\texttt{httpd.conf} configuration syntax in the following example:

\begin{verbatim}
<Location /svn>
  DAV svn

  # Automatically map any "/svn/foo" URL to repository /var/svn/foo
  SVNParentPath /var/svn
</Location>
\end{verbatim}

Using this syntax, Apache will delegate the handling of all URLs whose
path portions begin with \texttt{/svn/} to the Subversion DAV provider,
which will then assume that any items in the directory specified by the
\texttt{SVNParentPath} directive are actually Subversion repositories.
This is a particularly convenient syntax in that, unlike the use of the
\texttt{SVNPath} directive, you don\textquotesingle t have to restart
Apache to add or remove hosted repositories.

Be sure that when you define your new \texttt{Location}, it
doesn\textquotesingle t overlap with other exported locations. For
example, if your main \texttt{DocumentRoot} is exported to
\texttt{/www}, do not export a Subversion repository in
\texttt{\textless{}Location\ /www/repos\textgreater{}}. If a request
comes in for the URI \texttt{/www/repos/foo.c}, Apache
won\textquotesingle t know whether to look for a file
\texttt{repos/foo.c} in the \texttt{DocumentRoot}, or whether to
delegate \texttt{mod\_dav\_svn} to return \texttt{foo.c} from the
Subversion repository. The result is often an error from the server of
the form \texttt{301\ Moved\ Permanently}.

Server Names and the COPY Request

Subversion makes use of the \texttt{COPY} request type to perform
server-side copies of files and directories. As part of the sanity
checking done by the Apache modules, the source of the copy is expected
to be located on the same machine as the destination of the copy. To
satisfy this requirement, you might need to tell \texttt{mod\_dav} the
name you use as the hostname of your server. Generally, you can use the
\texttt{ServerName} directive in \texttt{httpd.conf} to accomplish this.

\begin{verbatim}
ServerName svn.example.com
\end{verbatim}

If you are using Apache\textquotesingle s virtual hosting support via
the \texttt{NameVirtualHost} directive, you may need to use the
\texttt{ServerAlias} directive to specify additional names by which your
server is known. Again, refer to the Apache documentation for full
details.

At this stage, you should strongly consider the question of permissions.
If you\textquotesingle ve been running Apache for some time now as your
regular web server, you probably already have a collection of
content---web pages, scripts, and such. These items have already been
configured with a set of permissions that allows them to work with
Apache, or more appropriately, that allows Apache to work with those
files. Apache, when used as a Subversion server, will also need the
correct permissions to read and write to your Subversion repository.

You will need to determine a permission system setup that satisfies
Subversion\textquotesingle s requirements without messing up any
previously existing web page or script installations. This might mean
changing the permissions on your Subversion repository to match those in
use by other things that Apache serves for you, or it could mean using
the \texttt{User} and \texttt{Group} directives in \texttt{httpd.conf}
to specify that Apache should run as the user and group that owns your
Subversion repository. There is no single correct way to set up your
permissions, and each administrator will have different reasons for
doing things a certain way. Just be aware that permission-related
problems are perhaps the most common oversight when configuring a
Subversion repository for use with Apache.

\subsection{Authentication Options}\label{svn.serverconfig.httpd.authn}

At this point, if you configured \texttt{httpd.conf} to contain
something such as the following:

\begin{verbatim}
<Location /svn>
  DAV svn
  SVNParentPath /var/svn
</Location>
\end{verbatim}

your repository is ``anonymously'' accessible to the world. Until you
configure some authentication and authorization policies, the Subversion
repositories that you make available via the \texttt{Location} directive
will be generally accessible to everyone. In other words:

\begin{itemize}
\item
  Anyone can use a Subversion client to check out a working copy of a
  repository URL (or any of its subdirectories).
\item
  Anyone can interactively browse the repository\textquotesingle s
  latest revision simply by pointing a web browser to the repository
  URL.
\item
  Anyone can commit to the repository.
\end{itemize}

Of course, you might have already set up a pre-commit hook script to
prevent commits (see \hyperref[svn.reposadmin.hooks]{Implementing
Repository Hooks}). But as you read on, you\textquotesingle ll see that
it\textquotesingle s also possible to use Apache\textquotesingle s
built-in methods to restrict access in specific ways.

Requiring authentication defends against invalid users directly
accessing the repository, but does not guard the privacy of valid
users\textquotesingle{} network activity. See
\hyperref[svn.serverconfig.httpd.ssl]{Protecting network traffic with
SSL} for how to configure your server to support SSL encryption, which
can provide that extra layer of protection.

\subsubsection{Basic
authentication}\label{svn.serverconfig.httpd.authn.basic}

The easiest way to authenticate a client is via the HTTP Basic
authentication mechanism, which simply uses a username and password to
verify a user\textquotesingle s identity. Apache provides the
\texttt{htpasswd} utility\footnote{See
  \href{https://httpd.apache.org/docs/current/programs/htpasswd.html}{}.}
for managing files containing usernames and passwords.

Basic authentication is \emph{extremely} insecure, because it sends
passwords over the network in very nearly plain text. See
\hyperref[svn.serverconfig.httpd.authn.digest]{Digest authentication}
for details on using the much safer Digest mechanism.

First, create a password file and grant access to users Harry and Sally:

\begin{verbatim}
$ ### First time: use -c to create the file
$ ### Use -m to use MD5 encryption of the password, which is more secure
$ htpasswd -c -m /etc/svn-auth.htpasswd harry
New password: *****
Re-type new password: *****
Adding password for user harry
$ htpasswd -m /etc/svn-auth.htpasswd sally
New password: *******
Re-type new password: *******
Adding password for user sally
$
\end{verbatim}

Next, ensure that Apache has access to the modules which provide the
Basic authentication and related functionality:
\texttt{mod\_auth\_basic}, \texttt{mod\_authn\_file}, and
\texttt{mod\_authz\_user}. In many cases, these modules are compiled
into \texttt{httpd} itself, but if not, you might need to explicitly
load one or more of them using the \texttt{LoadModule} directive:

\begin{verbatim}
LoadModule auth_basic_module   modules/mod_auth_basic.so
LoadModule authn_file_module   modules/mod_authn_file.so
LoadModule authz_user_module   moduels/mod_authz_user.so
\end{verbatim}

After ensuring the Apache has access to the required functionality,
you\textquotesingle ll need to add some more directives inside the
\texttt{\textless{}Location\textgreater{}} block to tell Apache what
type of authentication you wish to use, and just how to do so:

\begin{verbatim}
<Location /svn>
  DAV svn
  SVNParentPath /var/svn

  # Authentication: Basic
  AuthName "Subversion repository"
  AuthType Basic
  AuthBasicProvider file
  AuthUserFile /etc/svn-auth.htpasswd
</Location>
\end{verbatim}

These directives work as follows:

\begin{itemize}
\item
  \texttt{AuthName} is an arbitrary name that you choose for the
  authentication domain. Most browsers display this name in the dialog
  box when prompting for username and password.
\item
  \texttt{AuthType} specifies the type of authentication to use.
\item
  \texttt{AuthBasicProvider} specifies the Basic authentication provider
  to use for the location. In our example, we wish to consult a local
  password file.
\item
  \texttt{AuthUserFile} specifies the location of the password file to
  use.
\end{itemize}

However, this \texttt{\textless{}Location\textgreater{}} block
doesn\textquotesingle t yet do anything useful. It merely tells Apache
that \emph{if} authorization were required, it should challenge the
Subversion client for a username and password. (When authorization is
required, Apache requires authentication as well.)
What\textquotesingle s missing here, however, are directives that tell
Apache \emph{which sorts} of client requests require authorization;
currently, none do. The simplest thing is to specify that \emph{all}
requests require authorization by adding \texttt{Require\ valid-user} to
the block:

\begin{verbatim}
<Location /svn>
  DAV svn
  SVNParentPath /var/svn

  # Authentication: Basic
  AuthName "Subversion repository"
  AuthType Basic
  AuthBasicProvider file
  AuthUserFile /etc/svn-auth.htpasswd

  # Authorization: Authenticated users only
  Require valid-user
</Location>
\end{verbatim}

Refer to \hyperref[svn.serverconfig.httpd.authz]{Authorization Options}
for more detail on the \texttt{Require} directive and other ways to set
authorization policies.

The default value of the \texttt{AuthBasicProvider} option is
\texttt{file}, so we won\textquotesingle t bother including it in future
examples. Just know that if in a broader context you\textquotesingle ve
set this value to something else, you\textquotesingle ll need to
explicitly reset it to \texttt{file} within your Subversion
\texttt{\textless{}Location\textgreater{}} block in order to get that
behavior.

\subsubsection{Digest
authentication}\label{svn.serverconfig.httpd.authn.digest}

Digest authentication is an improvement on Basic authentication which
allows the server to verify a client\textquotesingle s identity without
sending the password over the network unprotected. Both client and
server create a non-reversible MD5 hash of the username, password,
requested URI, and a nonce (number used once) provided by the server and
changed each time authentication is required. The client sends its hash
to the server, and the server then verifies that the hashes match.

Configuring Apache to use Digest authentication is straightforward.
You\textquotesingle ll need to ensure that the
\texttt{mod\_auth\_digest} module is available (instead of
\texttt{mod\_auth\_basic}), and then make a few small variations on our
prior example:

\begin{verbatim}
<Location /svn>
  DAV svn
  SVNParentPath /var/svn

  # Authentication: Digest
  AuthName "Subversion repository"
  AuthType Digest
  AuthDigestProvider file
  AuthUserFile /etc/svn-auth.htdigest

  # Authorization: Authenticated users only
  Require valid-user
</Location>
\end{verbatim}

Notice that \texttt{AuthType} is now set to \texttt{Digest}, and we
specify a different path for \texttt{AuthUserFile}. Digest
authentication uses a different file format than Basic authentication,
created and managed using Apache\textquotesingle s \texttt{htdigest}
utility\footnote{See
  \href{https://httpd.apache.org/docs/current/programs/htdigest.html}{}.}
rather than \texttt{htpasswd}. Digest authentication also has the
additional concept of a ``realm'', which must match the value of the
\texttt{AuthName} directive.

For Digest authentication, the authentication provider is selected using
the \texttt{AuthDigestProvider} as shown in the previous example. As was
the case with the \texttt{AuthBasicProvider} directive, \texttt{file} is
the default value of the \texttt{AuthDigestProvider} option, so this
line is not strictly required unless you need to override a different
value thereof inherited from a broader configuration context.

The password file can be created as follows:

\begin{verbatim}
$ ### First time: use -c to create the file
$ htdigest -c /etc/svn-auth.htdigest "Subversion repository" harry
Adding password for harry in realm Subversion repository.
New password: *****
Re-type new password: *****
$ htdigest /etc/svn-auth.htdigest "Subversion repository" sally
Adding user sally in realm Subversion repository
New password: *******
Re-type new password: *******
$
\end{verbatim}

\subsection{Authorization Options}\label{svn.serverconfig.httpd.authz}

At this point, you\textquotesingle ve configured authentication, but not
authorization. Apache is able to challenge clients and confirm
identities, but it has not been told how to allow or restrict access to
the clients bearing those identities. This section describes two
strategies for controlling access to your repositories.

\subsubsection{Blanket access
control}\label{svn.serverconfig.httpd.authz.blanket}

The simplest form of access control is to authorize certain users for
either read-only access to a repository or read/write access to a
repository.

You can restrict access on all repository operations by adding
\texttt{Require\ valid-user} directly inside the
\texttt{\textless{}Location\textgreater{}} block. The example from
\hyperref[svn.serverconfig.httpd.authn.digest]{Digest authentication}
allows only clients that successfully authenticate to do anything with
the Subversion repository:

\begin{verbatim}
<Location /svn>
  DAV svn
  SVNParentPath /var/svn

  # Authentication: Digest
  AuthName "Subversion repository"
  AuthType Digest
  AuthUserFile /etc/svn-auth.htdigest

  # Authorization: Authenticated users only
  Require valid-user
</Location>
\end{verbatim}

Sometimes you don\textquotesingle t need to run such a tight ship. For
example, the server hosting Subversion\textquotesingle s own source code
at \href{https://svn.apache.org/repos/asf/subversion/}{} allows anyone
in the world to perform read-only repository tasks (such as checking out
working copies and browsing the repository), but restricts write
operations to authenticated users. The \texttt{Limit} and
\texttt{LimitExcept} directives allow for this type of selective
restriction. Like the \texttt{Location} directive, these blocks have
starting and ending tags, and you would nest them inside your
\texttt{\textless{}Location\textgreater{}} block.

The parameters present on the \texttt{Limit} and \texttt{LimitExcept}
directives are HTTP request types that are affected by that block. For
example, to allow anonymous read-only operations, you would use the
\texttt{LimitExcept} directive (passing the \texttt{GET},
\texttt{PROPFIND}, \texttt{OPTIONS}, and \texttt{REPORT} request type
parameters) and place the previously mentioned
\texttt{Require\ valid-user} directive inside the
\texttt{\textless{}LimitExcept\textgreater{}} block instead of just
inside the \texttt{\textless{}Location\textgreater{}} block.

\begin{verbatim}
<Location /svn>
  DAV svn
  SVNParentPath /var/svn

  # Authentication: Digest
  AuthName "Subversion repository"
  AuthType Digest
  AuthUserFile /etc/svn-auth.htdigest

  # Authorization: Authenticated users only for non-read-only
  #                (write) operations; allow anonymous reads
  <LimitExcept GET PROPFIND OPTIONS REPORT>
    Require valid-user
  </LimitExcept>
</Location>
\end{verbatim}

These are only a few simple examples. For more in-depth information
about Apache access control and the \texttt{Require} directive, take a
look at the \texttt{Security} section of the Apache
documentation\textquotesingle s tutorials collection at
\href{https://httpd.apache.org/docs-2.0/misc/tutorials.html}{}.

\subsubsection{Per-directory access
control}\label{svn.serverconfig.httpd.authz.perdir}

It\textquotesingle s possible to set up finer-grained permissions using
\texttt{mod\_authz\_svn}. This Apache module grabs the various opaque
URLs passing from client to server, asks \texttt{mod\_dav\_svn} to
decode them, and then possibly vetoes requests based on access policies
defined in a configuration file.

If you\textquotesingle ve built Subversion from source code,
\texttt{mod\_authz\_svn} is automatically built and installed alongside
\texttt{mod\_dav\_svn}. Many binary distributions install it
automatically as well. To verify that it\textquotesingle s installed
correctly, make sure it comes right after
\texttt{mod\_dav\_svn}\textquotesingle s \texttt{LoadModule} directive
in \texttt{httpd.conf}:

\begin{verbatim}
LoadModule dav_module         modules/mod_dav.so
LoadModule dav_svn_module     modules/mod_dav_svn.so
LoadModule authz_svn_module   modules/mod_authz_svn.so
\end{verbatim}

To activate this module, you need to configure your
\texttt{\textless{}Location\textgreater{}} block to use the
\texttt{AuthzSVNAccessFile} directive which specifies a single file
containing the permissions policy for paths within your repositories.
Beginning with Subversion 1.7, you can also use
\texttt{AuthzSVNReposRelativeAccessFile} directive to specify a per
repository access file. (In a moment, we\textquotesingle ll discuss the
format of that file.)

Apache is flexible, so you have the option to configure your block in
one of three general patterns. To begin, choose one of these basic
configuration patterns. (The following examples are very simple; look at
Apache\textquotesingle s own documentation for much more detail on
Apache authentication and authorization options.)

The most open approach is to allow access to everyone. This means Apache
never sends authentication challenges, and all users are treated as
``anonymous''. (See
\hyperref[svn.serverconfig.httpd.authz.perdir.ex-1]{A sample
configuration for anonymous access}.)

\protect\phantomsection\label{svn.serverconfig.httpd.authz.perdir.ex-1}
A sample configuration for anonymous access

\begin{verbatim}
<Location /repos>
  DAV svn
  SVNParentPath /var/svn

  # Authentication: None

  # Authorization: Path-based access control
  AuthzSVNAccessFile /path/to/access/file
</Location>
\end{verbatim}

On the opposite end of the paranoia scale, you can configure Apache to
authenticate all clients. This block unconditionally requires
authentication via the \texttt{Require\ valid-user} directive, and
defines a means to authenticate valid users. (See
\hyperref[svn.serverconfig.httpd.authz.perdir.ex-2]{A sample
configuration for authenticated access}.)

\protect\phantomsection\label{svn.serverconfig.httpd.authz.perdir.ex-2}
A sample configuration for authenticated access

\begin{verbatim}
<Location /repos>
  DAV svn
  SVNParentPath /var/svn

  # Authentication: Digest
  AuthName "Subversion repository"
  AuthType Digest
  AuthUserFile /etc/svn-auth.htdigest

  # Authorization: Path-based access control; authenticated users only
  AuthzSVNAccessFile /path/to/access/file
  Require valid-user
</Location>
\end{verbatim}

A third very popular pattern is to allow a combination of authenticated
and anonymous access. For example, many administrators want to allow
anonymous users to read certain repository directories, but restrict
access to more sensitive areas to authenticated users. In this setup,
all users start out accessing the repository anonymously. If your access
control policy demands a real username at any point, Apache will demand
authentication from the client. To do this, use both the
\texttt{Satisfy\ Any} and \texttt{Require\ valid-user} directives. (See
\hyperref[svn.serverconfig.httpd.authz.perdir.ex-3]{A sample
configuration for mixed authenticated/anonymous access}.)

\protect\phantomsection\label{svn.serverconfig.httpd.authz.perdir.ex-3}
A sample configuration for mixed authenticated/anonymous access

\begin{verbatim}
<Location /repos>
  DAV svn
  SVNParentPath /var/svn

  # Authentication: Digest
  AuthName "Subversion repository"
  AuthType Digest
  AuthUserFile /etc/svn-auth.htdigest

  # Authorization: Path-based access control; try anonymous access
  #                first, but authenticate if necessary
  AuthzSVNAccessFile /path/to/access/file
  Satisfy Any
  Require valid-user
</Location>
\end{verbatim}

The next step is to create the authorization file containing access
rules for particular paths within the repository. We describe how later
in this chapter, in
\hyperref[svn.serverconfig.pathbasedauthz]{Path-Based Authorization}.

\subsubsection{Disabling path-based
checks}\label{svn.serverconfig.httpd.authz.pathauthzoff}

The \texttt{mod\_dav\_svn} module goes through a lot of work to make
sure that data you\textquotesingle ve marked ``unreadable''
doesn\textquotesingle t get accidentally leaked. This means it needs to
closely monitor all of the paths and file-contents returned by commands
such as \texttt{svn\ checkout} and \texttt{svn\ update}. If these
commands encounter a path that isn\textquotesingle t readable according
to some authorization policy, the path is typically omitted altogether.
In the case of history or rename tracing---for example, running a
command such as \texttt{svn\ cat\ -r\ OLD\ foo.c} on a file that was
renamed long ago---the rename tracking will simply halt if one of the
object\textquotesingle s former names is determined to be
read-restricted.

All of this path checking can sometimes be quite expensive, especially
in the case of \texttt{svn\ log}. When retrieving a list of revisions,
the server looks at every changed path in each revision and checks it
for readability. If an unreadable path is discovered,
it\textquotesingle s omitted from the list of the
revision\textquotesingle s changed paths (normally seen with the
\texttt{-\/-verbose} (\texttt{-v}) option), and the whole log message is
suppressed. Needless to say, this can be time-consuming on revisions
that affect a large number of files. This is the cost of security: even
if you haven\textquotesingle t configured a module such as
\texttt{mod\_authz\_svn} at all, the \texttt{mod\_dav\_svn} module is
still asking Apache \texttt{httpd} to run authorization checks on every
path. The \texttt{mod\_dav\_svn} module has no idea what authorization
modules have been installed, so all it can do is ask Apache to invoke
whatever might be present.

On the other hand, there\textquotesingle s also an escape hatch of
sorts, which allows you to trade security features for speed. If
you\textquotesingle re not enforcing any sort of per-directory
authorization (i.e., not using \texttt{mod\_authz\_svn} or similar
module), you can disable all of this path checking. In your
\texttt{httpd.conf} file, use the \texttt{SVNPathAuthz} directive as
shown in
\hyperref[svn.serverconfig.httpd.authz.pathauthzoff.ex-1]{Disabling path
checks altogether}.

\protect\phantomsection\label{svn.serverconfig.httpd.authz.pathauthzoff.ex-1}
Disabling path checks altogether

\begin{verbatim}
<Location /repos>
  DAV svn
  SVNParentPath /var/svn

  SVNPathAuthz off
</Location>
\end{verbatim}

The \texttt{SVNPathAuthz} directive is ``on'' by default. When set to
``off,'' all path-based authorization checking is disabled;
\texttt{mod\_dav\_svn} stops invoking authorization checks on every path
it discovers.

\subsubsection{Versioned in repository access
files}\label{svn.serverconfig.httpd.authz.inrepo-authz}

Beginning with Subversion 1.8, access files can be stored inside a
Subversion repository. It is possible to store the access file in the
same repository to which the access rules are being applied, or you can
store it in an entirely different repository. This approach enables
versioning features of Subversion for the path-based authorization
configuration.

Both \texttt{AuthzSVNAccessFile} and
\texttt{AuthzSVNReposRelativeAccessFile} configuration directives allow
to specify in-repository access file\textquotesingle s location. The
directives accept absolute \texttt{file://} URLs and repository relative
URLs (one which begins with \texttt{\^{}/}).

For example, it is possible to specify an absolute URL to in-repository
access file as shown in
\hyperref[svn.serverconfig.httpd.authz.inrepo-authz.ex-1]{Using single
versioned in repo access file}.

\protect\phantomsection\label{svn.serverconfig.httpd.authz.inrepo-authz.ex-1}
Using single versioned in repo access file

\begin{verbatim}
<Location /repos>
  DAV svn
  SVNParentPath /var/svn
  AuthzSVNAccessFile file:///var/svn/authzrepo/authz
</Location>
\end{verbatim}

You can also specify a relative URL to an in repository access file as
demonstrated in
\hyperref[svn.serverconfig.httpd.authz.inrepo-authz.ex-2]{Using per
repository in repo access files}.

\protect\phantomsection\label{svn.serverconfig.httpd.authz.inrepo-authz.ex-2}
Using per repository in repo access files

\begin{verbatim}
<Location /repos>
  DAV svn
  SVNParentPath /var/svn
  AuthzSVNReposRelativeAccessFile ^/authz
</Location>
\end{verbatim}

\subsection{Protecting network traffic with
SSL}\label{svn.serverconfig.httpd.ssl}

Connecting to a repository via \texttt{http://} means that all
Subversion activity is sent across the network in the clear. This means
that actions such as checkouts, commits, and updates could potentially
be intercepted by an unauthorized party ``sniffing'' network traffic.
Encrypting traffic using SSL is a good way to protect potentially
sensitive information over the network.

If a Subversion client is compiled to use OpenSSL, it gains the ability
to speak to an Apache server via \texttt{https://} URLs, so all traffic
is encrypted with a per-connection session key. The WebDAV library used
by the Subversion client is not only able to verify server certificates,
but can also supply client certificates when challenged by the server.

\subsubsection{Subversion server SSL certificate
configuration}\label{svn.serverconfig.httpd.ssl.server}

It\textquotesingle s beyond the scope of this book to describe how to
generate client and server SSL certificates and how to configure Apache
to use them. Many other references, including Apache\textquotesingle s
own documentation (\href{https://httpd.apache.org/docs/current/ssl/}{}),
describe the process.

SSL certificates from well-known entities generally cost money, but at a
bare minimum, you can configure Apache to use a self-signed certificate
generated with a tool such as OpenSSL
(\href{https://www.openssl.org}{}).\footnote{While self-signed
  certificates are still vulnerable to a ``man-in-the-middle'' attack
  (before a client sees the certificate for the first time), such an
  attack is much more difficult for a casual observer to pull off,
  compared to sniffing unprotected passwords.}

\subsubsection{Subversion client SSL certificate
management}\label{svn.serverconfig.httpd.ssl.client}

When connecting to Apache via \texttt{https://}, a Subversion client can
receive two different types of responses:

\begin{itemize}
\item
  A server certificate
\item
  A challenge for a client certificate
\end{itemize}

\paragraph{Server
certificate}\label{svn.serverconfig.httpd.ssl.client.servercert}

When the client receives a server certificate, it needs to verify that
the server is who it claims to be. OpenSSL does this by examining the
signer of the server certificate, or certificate authority (CA). If
OpenSSL is unable to automatically trust the CA, or if some other
problem occurs (such as an expired certificate or hostname mismatch),
the Subversion command-line client will ask you whether you want to
trust the server certificate anyway:

\begin{verbatim}
$ svn list https://host.example.com/repos/project

Error validating server certificate for 'https://host.example.com:443':
 - The certificate is not issued by a trusted authority.  Use the
   fingerprint to validate the certificate manually!
Certificate information:
 - Hostname: host.example.com
 - Valid: from Jan 30 19:23:56 2004 GMT until Jan 30 19:23:56 2006 GMT
 - Issuer: CA, example.com, Sometown, California, US
 - Fingerprint: 7d:e1:a9:34:33:39:ba:6a:e9:a5:c4:22:98:7b:76:5c:92:a0:9c:7b

(R)eject, accept (t)emporarily or accept (p)ermanently?
\end{verbatim}

This dialogue is essentially the same question you may have seen coming
from your web browser (which is just another HTTP client like
Subversion). If you choose the \texttt{(p)}ermanent option, Subversion
will cache the server certificate in your private runtime \texttt{auth/}
area, just as your username and password are cached (see
\hyperref[svn.serverconfig.netmodel.credcache]{Caching credentials}),
and will automatically trust the certificate in the future.

Your runtime \texttt{servers} file also gives you the ability to make
your Subversion client automatically trust specific CAs, either globally
or on a per-host basis. Simply set the \texttt{ssl-authority-files}
variable to a semicolon-separated list of PEM-encoded CA certificates:

\begin{verbatim}
[global]
ssl-authority-files = /path/to/CAcert1.pem;/path/to/CAcert2.pem
\end{verbatim}

Many OpenSSL installations also have a predefined set of ``default'' CAs
that are nearly universally trusted. To make the Subversion client
automatically trust these standard authorities, set the
\texttt{ssl-trust-default-ca} variable to \texttt{true}.

\paragraph{Client certificate
challenge}\label{svn.serverconfig.httpd.ssl.client.clientcert}

If the client receives a challenge for a certificate, the server is
asking the client to prove its identity. The client must send back a
certificate signed by a CA that the server trusts, along with a
challenge response which proves that the client owns the private key
associated with the certificate. The private key and certificate are
usually stored in an encrypted format on disk, protected by a
passphrase. When Subversion receives this challenge, it will ask you for
the path to the encrypted file and the passphrase that protects it:

\begin{verbatim}
$ svn list https://host.example.com/repos/project

Authentication realm: https://host.example.com:443
Client certificate filename: /path/to/my/cert.p12
Passphrase for '/path/to/my/cert.p12':  ********
\end{verbatim}

Notice that the client credentials are stored in a \texttt{.p12} file.
To use a client certificate with Subversion, it must be in PKCS\#12
format, which is a portable standard. Most web browsers are able to
import and export certificates in that format. Another option is to use
the OpenSSL command-line tools to convert existing certificates into
PKCS\#12.

The runtime \texttt{servers} file also allows you to automate this
challenge on a per-host basis. If you set the
\texttt{ssl-client-cert-file} and \texttt{ssl-client-cert-password}
variables, Subversion can automatically respond to a client certificate
challenge without prompting you:

\begin{verbatim}
[groups]
examplehost = host.example.com

[examplehost]
ssl-client-cert-file = /path/to/my/cert.p12
ssl-client-cert-password = somepassword
\end{verbatim}

More security-conscious folk might want to exclude
\texttt{ssl-client-cert-password} to avoid storing the passphrase in the
clear on disk.

\subsection{Tuning for Performance}\label{svn.serverconfig.httpd.perf}

The Apache HTTP Server is built for performance, but you can improve
upon its default configuration to get even better results out of your
Subversion service offering. In this section, we\textquotesingle ll
recommend some specific configuration changes to consider. Understand,
however, that some of the \texttt{httpd.conf} configuration options
we\textquotesingle ll be discussing herein affect the general behavior
of your server, not merely the Subversion service. As such, you need to
consider the full breadth of your HTTP service offering to discern how
modifications to these settings for Subversion\textquotesingle s sake
may affect your other services.

\subsubsection{KeepAlive}\label{svn.serverconfig.httpd.perf.keepalive}

By default, the Apache HTTP Server is configured to enable the re-use of
a single server connection for multiple requests. That\textquotesingle s
very beneficial for Subversion, because unlike many HTTP-based
applications, Subversion can very quickly generate hundreds or thousands
of requests against the server for a single operation, and the cost of
opening a new connection to the server is non-trivial. Subversion wants
to squeeze as many requests as possible out of a single connection
before that connection is terminated by the server. The
\texttt{KeepAlive} directive is the boolean flag which enables or
disables this connection re-use facility, and as we indicated
previously, by default its value is \texttt{On}.

But there\textquotesingle s another directive which limits the number of
requests a client is allowed to submit on a single connection: the
\texttt{MaxKeepAliveRequests} directive. The default value for that
option is \texttt{100}. This was probably sufficient for older versions
of Subversion, but Subversion 1.8 employs a different HTTP
communications library (called Serf) which prefers to pipeline several
smaller requests for specific bits of information rather than asking the
server to transmit huge chunks of data in a single response. We
recommend that you increase the value of the
\texttt{MaxKeepAliveRequests} option to at least \texttt{1000}.

\begin{verbatim}
#
# KeepAlive: Whether or not to allow persistent connections (more than
# one request per connection). Set to "Off" to deactivate.
#
KeepAlive On

#
# MaxKeepAliveRequests: The maximum number of requests to allow
# during a persistent connection. Set to 0 to allow an unlimited amount.
# We recommend you leave this number high, for maximum performance.
#
MaxKeepAliveRequests 1000
\end{verbatim}

\subsubsection{Bulk
updates}\label{svn.serverconfig.httpd.perf.bulk-updates}

The biggest difference between the way that Subversion 1.8 clients and
pre-1.8 clients behave is in how update-style operations
(\texttt{svn\ checkout}, \texttt{svn\ update}, \texttt{svn\ switch},
etc.) are handled. Older clients which used the Neon HTTP library for
communications preferred to ask the server for the entire payload of
information required from the server in a single request. Admins will
have noticed that in their server logs, there would be some initial
handshaking operations, and then a \texttt{REPORT} request with a
massive response. That response was the entire checkout/update dataset!

Subversion clients which use the Serf HTTP library---which includes all
clients built atop the Subversion 1.8---still send the \texttt{REPORT}
request, but with slightly different flags set inside that request.
These flags ask the server not to send all the data for the operation,
but to instead send only a checklist of other more specific things that
the client needs to subsequently fetch from the server in order to
complete that operation. In the server\textquotesingle s
\texttt{access\_log}, that \texttt{REPORT} is followed by many smaller
requests (\texttt{GET}s and, in older versions of Subversion,
\texttt{PROPFIND}s).

There are pros and cons to each approach. As we\textquotesingle ve
mentioned, the so-called bulk updates generate considerably less
information in the server logs, but a given Apache HTTP Server child
process is completely consumed for the duration of what could be a
lengthy operation. Non-bulk updates offer opportunities for setting up
content caches (which themselves can improve performance), but generate
server log traffic which is whole orders of magnitude larger than the
bulk update approach. So, for one reason or another, administrators may
desire to exert a little more control over which approach the clients
use. Subversion 1.6 introduced the \texttt{SVNAllowBulkUpdates}
\texttt{mod\_dav\_svn} directive---a simple boolean flag---to allow
admins to specify whether the server was allowed to honor bulk update
requests. In Subversion 1.8, this directive has expanded to include a
\texttt{Prefer} value in addition to the \texttt{On} and \texttt{Off}
values it already supported. When \texttt{SVNAllowBulkUpdates} is set to
\texttt{Prefer}, supporting clients (1.8 or newer) will try to use the
bulk update approach unless otherwise configured.

\subsection{Extra Goodies}\label{svn.serverconfig.httpd.extra}

We\textquotesingle ve covered most of the authentication and
authorization options for Apache and \texttt{mod\_dav\_svn}. But there
are a few other nice features that Apache provides.

\subsubsection{Repository
browsing}\label{svn.serverconfig.httpd.extra.browsing}

One of the most useful benefits of an Apache/WebDAV configuration for
your Subversion repository is that your versioned files and directories
are immediately available for viewing via a regular web browser. Since
Subversion uses URLs to identify versioned resources, those URLs used
for HTTP-based repository access can be typed directly into a web
browser. Your browser will issue an HTTP \texttt{GET} request for that
URL; based on whether that URL represents a versioned directory or file,
\texttt{mod\_dav\_svn} will respond with a directory listing or with
file contents.

\paragraph{URL
syntax}\label{svn.serverconfig.httpd.extra.browsing.syntax}

If the URLs do not contain any information about which version of the
resource you wish to see, \texttt{mod\_dav\_svn} will answer with the
youngest version. This functionality has the wonderful side effect that
you can pass around Subversion URLs to your peers as references to
documents, and those URLs will always point at the latest manifestation
of that document. Of course, you can even use the URLs as hyperlinks
from other web sites, too.

As of Subversion 1.6, \texttt{mod\_dav\_svn} supports a public URI
syntax for examining older revisions of both files and directories. The
syntax uses the query string portion of the URL to specify either or
both of a peg revision and operative revision, which Subversion will
then use to determine which version of the file or directory to display
to your web browser. Add the query string name/value pair
\texttt{p=PEGREV}, where \textless PEGREV\textgreater{} is a revision
number, to specify the peg revision you wish to apply to the request.
Use \texttt{r=REV}, where \textless REV\textgreater{} is a revision
number, to specify an operative revision.

For example, if you wish to see the latest version of a
\texttt{README.txt} file located in your project\textquotesingle s
\texttt{/trunk}, point your web browser to that file\textquotesingle s
repository URL, which might look something like the following:

\begin{verbatim}
http://host.example.com/repos/project/trunk/README.txt
\end{verbatim}

If you now wish to see some older version of that file, add an operative
revision to the URL\textquotesingle s query string:

\begin{verbatim}
http://host.example.com/repos/project/trunk/README.txt?r=1234
\end{verbatim}

What if the thing you\textquotesingle re trying to view no longer exists
in the youngest revision of the repository? That\textquotesingle s where
a peg revision is handy:

\begin{verbatim}
http://host.example.com/repos/project/trunk/deleted-thing.txt?p=321
\end{verbatim}

And of course, you can combine peg revision and operative revision
specifiers to fine-tune the exact item you wish to view:

\begin{verbatim}
http://host.example.com/repos/project/trunk/renamed-thing.txt?p=123&r=21
\end{verbatim}

The previous URL would display revision 21 of the object which, in
revision 123, was located at \texttt{/trunk/renamed-thing.txt} in the
repository. See \hyperref[svn.advanced.pegrevs]{Peg and Operative
Revisions} for a detailed explanation of these ``peg revision'' and
``operative revision'' concepts. They can be a bit tricky to wrap your
head around.

Beginning with Subversion 1.8, \texttt{mod\_dav\_svn} has the ability to
substitute keywords. When \texttt{mod\_dav\_svn} finds the query
argument \texttt{kw=1} added to the URL of a file, it will expand
keywords when delivering the file\textquotesingle s content. Omitting
the \texttt{kw} parameter or using any value besides \texttt{1} for that
parameter will cause Subversion to use its default behavior, which is to
deliver the file\textquotesingle s content without any keywords
expanded.

Because keyword substitution is typically performed by the Subversion
client as part of its working copy administration and management, this
is a handy way to get the server to deliver a keyword-expanded form of
your versioned file without the use of a working copy.

For example, if you wish to see the latest version of a
\texttt{README.txt} file located in your project\textquotesingle s
\texttt{/trunk} with keywords expanded, add the query argument
\texttt{kw=1} to the URL:

\begin{verbatim}
http://host.example.com/repos/project/trunk/README.txt?kw=1
\end{verbatim}

As with client-side keyword expansion, only those keywords which are
explicitly designated for expansion via the \texttt{svn:keywords}
property set on the file itself will be expanded. See
\hyperref[svn.advanced.props.special.keywords]{Keyword Substitution} for
a detailed description of the keyword substitution feature.

As a reminder, \texttt{mod\_dav\_svn} offers only a limited repository
browsing experience. You can see directory listings and file contents,
but no revision properties (such as commit log messages) or
file/directory properties. For folks who require more extensive browsing
of repositories and their history, there are several third-party
software packages which offer this. Some examples include ViewVC
(\href{https://viewvc.org}{}), Trac (\href{https://trac.edgewall.org}{})
and WebSVN (\href{https://websvnphp.github.io}{}). These third-party
tools don\textquotesingle t affect
\texttt{mod\_dav\_svn}\textquotesingle s built-in ``browseability'', and
generally offer a much wider set of features, including the display of
the aforementioned property sets, display of content differences between
file revisions, and so on.

\paragraph{Proper MIME
type}\label{svn.serverconfig.httpd.extra.browsing.mimetype}

When browsing a Subversion repository, the web browser gets a clue about
how to render a file\textquotesingle s contents by looking at the
\texttt{Content-Type:} header returned in Apache\textquotesingle s
response to the HTTP \texttt{GET} request. The value of this header is
some sort of MIME type. By default, Apache will tell the web browsers
that all repository files are of the ``default'' MIME type, typically
\texttt{text/plain}. This can be frustrating, however, if a user wishes
repository files to render as something more meaningful---for example,
it might be nice to have a \texttt{foo.html} file in the repository
actually render as HTML when browsing.

To make this happen, you need only to make sure that your files have the
proper \texttt{svn:mime-type} set. We discuss this in more detail in
\hyperref[svn.advanced.props.special.mime-type]{File Content Type}, and
you can even configure your client to automatically attach proper
\texttt{svn:mime-type} properties to files entering the repository for
the first time; see \hyperref[svn.advanced.props.auto]{Automatic
Property Setting}.

Continuing our example, if one were to set the \texttt{svn:mime-type}
property to \texttt{text/html} on file \texttt{foo.html}, Apache would
properly tell your web browser to render the file as HTML. One could
also attach proper \texttt{image/*} MIME-type properties to image files
and ultimately get an entire web site to be viewable directly from a
repository! There\textquotesingle s generally no problem with this, as
long as the web site doesn\textquotesingle t contain any dynamically
generated content.

\paragraph{Customizing the
look}\label{svn.serverconfig.httpd.extra.browsing.xslt}

You generally will get more use out of URLs to versioned files---after
all, that\textquotesingle s where the interesting content tends to lie.
But you might have occasion to browse a Subversion directory listing,
where you\textquotesingle ll quickly note that the generated HTML used
to display that listing is very basic, and certainly not intended to be
aesthetically pleasing (or even interesting). To enable customization of
these directory displays, Subversion provides an XML index feature. A
single \texttt{SVNIndexXSLT} directive in your
repository\textquotesingle s \texttt{Location} block of
\texttt{httpd.conf} will instruct \texttt{mod\_dav\_svn} to generate XML
output when displaying a directory listing, and to reference the XSLT
stylesheet of your choice:

\begin{verbatim}
<Location /svn>
  DAV svn
  SVNParentPath /var/svn
  SVNIndexXSLT "/svnindex.xsl"
  …
</Location>
\end{verbatim}

Using the \texttt{SVNIndexXSLT} directive and a creative XSLT
stylesheet, you can make your directory listings match the color schemes
and imagery used in other parts of your web site. Or, if
you\textquotesingle d prefer, you can use the sample stylesheets
provided in the Subversion source distribution\textquotesingle s
\texttt{tools/xslt/} directory. Keep in mind that the path provided to
the \texttt{SVNIndexXSLT} directive is actually a URL path---browsers
need to be able to read your stylesheets to make use of them!

\paragraph{Listing
repositories}\label{svn.serverconfig.httpd.extra.browsing.reposlisting}

If you\textquotesingle re serving a collection of repositories from a
single URL via the \texttt{SVNParentPath} directive, then
it\textquotesingle s also possible to have Apache display all available
repositories to a web browser. Just activate the
\texttt{SVNListParentPath} directive:

\begin{verbatim}
<Location /svn>
  DAV svn
  SVNParentPath /var/svn
  SVNListParentPath on
  …
</Location>
\end{verbatim}

If a user now points her web browser to the URL
\texttt{http://host.example.com/svn/}, she\textquotesingle ll see a list
of all Subversion repositories sitting in \texttt{/var/svn}. Obviously,
this can be a security problem, so this feature is turned off by
default.

\subsubsection{Apache
logging}\label{svn.serverconfig.httpd.extra.logging}

Because Apache is an HTTP server at heart, it contains fantastically
flexible logging features. It\textquotesingle s beyond the scope of this
book to discuss all of the ways logging can be configured, but we should
point out that even the most generic \texttt{httpd.conf} file will cause
Apache to produce two logs: \texttt{error\_log} and
\texttt{access\_log}. These logs may appear in different places, but are
typically created in the logging area of your Apache installation. (On
Unix, they often live in \texttt{/usr/local/apache2/logs/}.)

The \texttt{error\_log} describes any internal errors that Apache runs
into as it works. The \texttt{access\_log} file records every incoming
HTTP request received by Apache. This makes it easy to see, for example,
which IP addresses Subversion clients are coming from, how often
particular clients use the server, which users are authenticating
properly, and which requests succeed or fail.

Unfortunately, because HTTP is a stateless protocol, even the simplest
Subversion client operation generates multiple network requests.
It\textquotesingle s very difficult to look at the \texttt{access\_log}
and deduce what the client was doing---most operations look like a
series of cryptic \texttt{PROPPATCH}, \texttt{GET}, \texttt{PUT}, and
\texttt{REPORT} requests. To make things worse, many client operations
send nearly identical series of requests, so it\textquotesingle s even
harder to tell them apart.

\texttt{mod\_dav\_svn}, however, can come to your aid. By activating an
``operational logging'' feature, you can ask \texttt{mod\_dav\_svn} to
create a separate log file describing what sort of high-level operations
your clients are performing.

To do this, you need to make use of Apache\textquotesingle s
\texttt{CustomLog} directive (which is explained in more detail in
Apache\textquotesingle s own documentation). Be sure to invoke this
directive \emph{outside} your Subversion \texttt{Location} block:

\begin{verbatim}
<Location /svn>
  DAV svn
  …
</Location>

CustomLog logs/svn_logfile "%t %u %{SVN-ACTION}e" env=SVN-ACTION
\end{verbatim}

In this example, we\textquotesingle re asking Apache to create a special
logfile, \texttt{svn\_logfile}, in the standard Apache \texttt{logs}
directory. The \texttt{\%t} and \texttt{\%u} variables are replaced by
the time and username of the request, respectively. The really important
parts are the two instances of \texttt{SVN-ACTION}. When Apache sees
that variable, it substitutes the value of the \texttt{SVN-ACTION}
environment variable, which is automatically set by
\texttt{mod\_dav\_svn} whenever it detects a high-level client action.

So, instead of having to interpret a traditional \texttt{access\_log}
like this:

\begin{verbatim}
[26/Jan/2007:22:25:29 -0600] "PROPFIND /svn/calc/!svn/vcc/default HTTP/1.1" 207 398
[26/Jan/2007:22:25:29 -0600] "PROPFIND /svn/calc/!svn/bln/59 HTTP/1.1" 207 449
[26/Jan/2007:22:25:29 -0600] "PROPFIND /svn/calc HTTP/1.1" 207 647
[26/Jan/2007:22:25:29 -0600] "REPORT /svn/calc/!svn/vcc/default HTTP/1.1" 200 607
[26/Jan/2007:22:25:31 -0600] "OPTIONS /svn/calc HTTP/1.1" 200 188
[26/Jan/2007:22:25:31 -0600] "MKACTIVITY /svn/calc/!svn/act/e6035ef7-5df0-4ac0-b811-4be7c823f998 HTTP/1.1" 201 227
…
\end{verbatim}

you can peruse a much more intelligible \texttt{svn\_logfile} like this:

\begin{verbatim}
[26/Jan/2007:22:24:20 -0600] - get-dir /tags r1729 props
[26/Jan/2007:22:24:27 -0600] - update /trunk r1729 depth=infinity
[26/Jan/2007:22:25:29 -0600] - status /trunk/foo r1729 depth=infinity
[26/Jan/2007:22:25:31 -0600] sally commit r1730
\end{verbatim}

In addition to the \texttt{SVN-ACTION} environment variable,
\texttt{mod\_dav\_svn} also populates the \texttt{SVN-REPOS} and
\texttt{SVN-REPOS-NAME} variables, which carry the filesystem path to
the repository and the basename thereof, respectively. You might wish to
include references to one or both of these variables in your
\texttt{CustomLog} format string, too, especially if you are combining
usage information from multiple repositories into a single log file. For
an exhaustive list of all actions logged, see
\hyperref[svn.serverconfig.operational-logging]{High-level Logging}.

Obviously, the more information that Apache logs about your Subversion
activities, the more disk space on your server those logs consume. It is
non uncommon for high-traffic Subversion servers to generate many
gigabytes of log information daily. Obviously, logs are only valuable if
they can be meaningfully processed, and huge log files can quickly
become unwieldy. There are various standard approaches to Apache HTTP
Server log management which are outside the scope of this book.
Administrators are encouraged to use the log rotation and archival
approach which works best for them.

But what if Subversion is simply generating too much log information to
be useful? For example, in
\hyperref[svn.serverconfig.httpd.perf.bulk-updates]{Bulk updates}, we
mentioned that certain approaches that Subversion clients may take to
checkouts and other update-style operations can cause rapid growth of
your server logs as requests for individual pieces of the update data
set are individually logged (whereas in previous version of Subversion,
they might not have been). In this case, you might consider using some
Apache configuration magic to selectively silence some of that log
activity.

Apache HTTP Server\textquotesingle s \texttt{mod\_setenvif} module
offers a \texttt{SetEnvIf} directive which is handy for conditionally
setting environment variables. And as it turns out, the
\texttt{CustomLog} directive can be told to conditionally log requests
based on the state of environment variables. The following is a sample
configuration which instructs the server to \emph{not} log \texttt{GET}
and \texttt{PROPFIND} requests aimed at private Subversion URLs.

\begin{verbatim}
# Matches everything, just to initialize the "in_repos" variable.
SetEnvIf Request_URI "^" in_repos=0

# Set "in_repos" if this is a request for a private Subversion URL.
SetEnvIf Request_URI "/!svn/" in_repos=1

# Set "do_not_log" for non-public request types we don't care to log.
SetEnvIf Request_Method "GET" do_not_log
SetEnvIf Request_Method "PROPFIND" do_not_log

# Unset "do_not_log" for URLs that aren't private Subversion URLs.
SetEnvIf in_repos 0 !do_not_log

# Log requests, but only if "do_not_log" isn't set.
CustomLog logs/access_log env=!do_not_log
\end{verbatim}

Using this configuration, \texttt{httpd} would still log \texttt{GET}
requests aimed at public Subversion URLs. These are the sorts of
requests generated by a web browser as someone navigates the repository
directly. But \texttt{GET} and \texttt{PROPFIND} requests aimed at
so-called ``private'' Subversion URLs---which are the very sorts of
requests used to fetch each and every individual file during a checkout
operation---won\textquotesingle t get logged.

\subsubsection{Write-through
proxying}\label{svn.serverconfig.httpd.extra.writethruproxy}

One of the nice advantages of using Apache as a Subversion server is
that it can be set up for simple replication. For example, suppose that
your team is distributed across four offices around the globe. The
Subversion repository can exist only in one of those offices, which
means the other three offices will not enjoy accessing
it---they\textquotesingle re likely to experience significantly slower
traffic and response times when updating and committing code. A powerful
solution is to set up a system consisting of one master Apache server
and several slave Apache servers. If you place a slave server in each
office, users can check out a working copy from whichever slave is
closest to them. All read requests go to their local slave. Write
requests get automatically routed to the single master server. When the
commit completes, the master then automatically ``pushes'' the new
revision to each slave server using the \texttt{svnsync} replication
tool.

This configuration creates a huge perceptual speed increase for your
users, because Subversion client traffic is typically 80--90\% read
requests. And if those requests are coming from a \emph{local} server,
it\textquotesingle s a huge win.

In this section, we\textquotesingle ll walk you through a standard setup
of this single-master/multiple-slave system. However, keep in mind that
your servers must be running at least Apache 2.2.0 (with
\texttt{mod\_proxy} loaded) and Subversion 1.5 (\texttt{mod\_dav\_svn}).

Ours is just one example of how you might setup a Subversion
write-through proxy configuration. There are other approaches. For
example, rather than having the master server push changes out to every
slave server, the slaves could periodically pull those changes from the
master. Or perhaps the master could push changes to a single slave,
which then pushes the same change to the next slave, and so on down the
line. Administrators are encouraged to use this section for basic
understanding of what happens in a Subversion WebDAV proxy deployment
scenario, and then implement the specific approach that works best for
their organization.

\paragraph{Configure the
servers}\label{svn.serverconfig.httpd.extra.writethruproxy.configure}

First, configure your master server\textquotesingle s
\texttt{httpd.conf} file in the usual way. Make the repository available
at a certain URI location, and configure authentication and
authorization however you\textquotesingle d like. After
that\textquotesingle s done, configure each of your ``slave'' servers in
the exact same way, but add the special \texttt{SVNMasterURI} directive
to the block:

\begin{verbatim}
<Location /svn>
  DAV svn
  SVNPath /var/svn/repos
  SVNMasterURI http://master.example.com/svn
  …
</Location>
\end{verbatim}

This new directive tells a slave server to redirect all write requests
to the master. (This is done automatically via Apache\textquotesingle s
\texttt{mod\_proxy} module.) Ordinary read requests, however, are still
serviced by the slaves. Be sure that your master and slave servers all
have matching authentication and authorization configurations; if they
fall out of sync, it can lead to big headaches.

Next, we need to deal with the problem of infinite recursion. With the
current configuration, imagine what will happen when a Subversion client
performs a commit to the master server. After the commit completes, the
server uses \texttt{svnsync} to replicate the new revision to each
slave. But because \texttt{svnsync} appears to be just another
Subversion client performing a commit, the slave will immediately
attempt to proxy the incoming write request back to the master! Hilarity
ensues.

The solution to this problem is to have the master push revisions to a
different \texttt{\textless{}Location\textgreater{}} on the slaves. This
location is configured to \emph{not} proxy write requests at all, but to
accept normal commits from (and only from) the master\textquotesingle s
IP address:

\begin{verbatim}
<Location /svn-proxy-sync>
  DAV svn
  SVNPath /var/svn/repos
  Order deny,allow
  Deny from all
  # Only let the server's IP address access this Location:
  Allow from 10.20.30.40
  …
</Location>
\end{verbatim}

\paragraph{Set up
replication}\label{svn.serverconfig.httpd.extra.writethruproxy.replicate}

Now that you\textquotesingle ve configured your \texttt{Location} blocks
on master and slaves, you need to configure the master to replicate to
the slaves. Our walkthough uses \texttt{svnsync}, which is covered in
more detail in
\hyperref[svn.reposadmin.maint.replication.svnsync]{Replication with
svnsync}.

First, make sure that each slave repository has a pre-revprop-change
hook script which allows remote revision property changes. (This is
standard procedure for being on the receiving end of \texttt{svnsync}.)
Then log into the master server and configure each of the slave
repository URIs to receive data from the master repository on the local
disk:

\begin{verbatim}
$ svnsync init http://slave1.example.com/svn-proxy-sync \
               file:///var/svn/repos
Copied properties for revision 0.
$ svnsync init http://slave2.example.com/svn-proxy-sync \
               file:///var/svn/repos
Copied properties for revision 0.
$ svnsync init http://slave3.example.com/svn-proxy-sync \
               file:///var/svn/repos
Copied properties for revision 0.

# Perform the initial replication

$ svnsync sync http://slave1.example.com/svn-proxy-sync \
               file:///var/svn/repos
Transmitting file data ....
Committed revision 1.
Copied properties for revision 1.
Transmitting file data .......
Committed revision 2.
Copied properties for revision 2.
…

$ svnsync sync http://slave2.example.com/svn-proxy-sync \
               file:///var/svn/repos
Transmitting file data ....
Committed revision 1.
Copied properties for revision 1.
Transmitting file data .......
Committed revision 2.
Copied properties for revision 2.
…

$ svnsync sync http://slave3.example.com/svn-proxy-sync \
               file:///var/svn/repos
Transmitting file data ....
Committed revision 1.
Copied properties for revision 1.
Transmitting file data .......
Committed revision 2.
Copied properties for revision 2.
…
\end{verbatim}

After this is done, we configure the master server\textquotesingle s
post-commit hook script to invoke \texttt{svnsync} on each slave server:

\begin{verbatim}
#!/bin/sh
# Post-commit script to replicate newly committed revision to slaves

svnsync sync http://slave1.example.com/svn-proxy-sync \
             file:///var/svn/repos > /dev/null 2>&1 &
svnsync sync http://slave2.example.com/svn-proxy-sync \
             file:///var/svn/repos > /dev/null 2>&1 &
svnsync sync http://slave3.example.com/svn-proxy-sync \
             file:///var/svn/repos > /dev/null 2>&1 &
\end{verbatim}

The extra bits on the end of each line aren\textquotesingle t necessary,
but they\textquotesingle re a sneaky way to allow the sync commands to
run in the background so that the Subversion client
isn\textquotesingle t left waiting forever for the commit to finish. In
addition to this post-commit hook, you\textquotesingle ll need a
post-revprop-change hook as well so that when a user, say, modifies a
log message, the slave servers get that change also:

\begin{verbatim}
#!/bin/sh
# Post-revprop-change script to replicate revprop-changes to slaves

REV=${2}
svnsync copy-revprops http://slave1.example.com/svn-proxy-sync \
                      file:///var/svn/repos \
                      -r ${REV} > /dev/null 2>&1 &
svnsync copy-revprops http://slave2.example.com/svn-proxy-sync \
                      file:///var/svn/repos \
                      -r ${REV} > /dev/null 2>&1 &
svnsync copy-revprops http://slave3.example.com/svn-proxy-sync \
                      file:///var/svn/repos \
                      -r ${REV} > /dev/null 2>&1 &
\end{verbatim}

The only thing we\textquotesingle ve left out here is what to do about
user-level locks (of the \texttt{svn\ lock} variety). Locks are enforced
by the master server during commit operations; but all the information
about locks is distributed during read operations such as
\texttt{svn\ update} and \texttt{svn\ status} by the slave server. As
such, a fully functioning proxy setup needs to perfectly replicate lock
information from the master server to the slave servers. Unfortunately,
most of the mechanisms that one might employ to accomplish this
replication fall short in one way or another\footnote{\href{https://subversion.apache.org/issue3457}{}
  tracks these problems.}. Many teams don\textquotesingle t use
Subversion\textquotesingle s locking features at all, so this may be a
nonissue for you. Sadly, for those teams which do use locks, we have no
recommendations on how to gracefully work around this shortcoming.

\paragraph{Caveats}\label{svn.serverconfig.httpd.extra.writethruproxy.caveats}

Your master/slave replication system should now be ready to use. A
couple of words of warning are in order, however. Remember that this
replication isn\textquotesingle t entirely robust in the face of
computer or network crashes. For example, if one of the automated
\texttt{svnsync} commands fails to complete for some reason, the slaves
will begin to fall behind. For example, your remote users will see that
they\textquotesingle ve committed revision 100, but then when they run
\texttt{svn\ update}, their local server will tell them that revision
100 doesn\textquotesingle t yet exist! Of course, the problem will be
automatically fixed the next time another commit happens and the
subsequent \texttt{svnsync} is successful---the sync will replicate all
waiting revisions. But still, you may want to set up some sort of
out-of-band monitoring to notice synchronization failures and force
\texttt{svnsync} to run when things go wrong.

Another limitation of the write-through proxy deployment model involves
version mismatches---of the version of Subversion which is installed,
that is---between the master and slave servers. Each new release of
Subversion may (and often does) add new features to the network protocol
used between the clients and servers. Since feature negotiation happens
against the slave, it is the slave\textquotesingle s protocol version
and feature set which is used. But write operations are passed through
to the master server quite literally. Therefore, there is a risk that
the slave server will answer a feature negotiation request from the
client in way that is true for the slave, but untrue for the master if
the master is running an older version of Subversion. This could result
in the client trying to use a new feature that the master
doesn\textquotesingle t understand, and failing.

Subversion 1.8 helps to mitigate this problem via the introduction of a
new Apache configuration directive, \texttt{SVNMasterVersion}. By
configuring each of your slave servers with \texttt{SVNMasterVersion}
set to the release version of the Subversion instance which is running
on your master server, the slave servers can more accurately negotiate
feature support with the client.

Unfortunately, Subversion 1.7 doesn\textquotesingle t offer the
\texttt{SVNMasterVersion} configuration directive and is known to have
some specific problems along these lines. If you are deploying a
Subversion 1.7 slave server in front of a pre-1.7 master,
you\textquotesingle ll want to configure your slave
server\textquotesingle s Subversion
\texttt{\textless{}Location\textgreater{}} block with the
\texttt{SVNAdvertiseV2Protocol\ Off} directive.

For the best results possible, try to run the same version of Subversion
on your master and slave servers.

Can We Set Up Replication with svnserve?

If you\textquotesingle re using \texttt{svnserve} instead of Apache as
your server, you can certainly configure your
repository\textquotesingle s hook scripts to invoke \texttt{svnsync} as
we\textquotesingle ve shown here, thereby causing automatic replication
from master to slaves. Unfortunately, at the time of this writing there
is no way to make slave \texttt{svnserve} servers automatically proxy
write requests back to the master server. This means your users would
only be able to check out read-only working copies from the slave
servers. You\textquotesingle d have to configure your slave servers to
disallow write access completely. This might be useful for creating
read-only ``mirrors'' of popular open source projects, but
it\textquotesingle s not a transparent proxying system.

\subsubsection{Other Apache
features}\label{svn.serverconfig.httpd.extra.other}

Several of the features already provided by Apache in its role as a
robust web server can be leveraged for increased functionality or
security in Subversion as well. The Subversion client is able to use SSL
(the Secure Sockets Layer, discussed earlier). If your Subversion client
is built to support SSL, it can access your Apache server using
\texttt{https://} and enjoy a high-quality encrypted network session.

Equally useful are other features of the Apache and Subversion
relationship, such as the ability to specify a custom port (instead of
the default HTTP port 80) or a virtual domain name by which the
Subversion repository should be accessed, or the ability to access the
repository through an HTTP proxy.

Finally, because \texttt{mod\_dav\_svn} is speaking a subset of the
WebDAV/DeltaV protocol, it\textquotesingle s possible to access the
repository via third-party DAV clients. Most modern operating systems
(Win32, OS X, and Linux) have the built-in ability to mount a DAV server
as a standard network ``shared folder.'' This is a complicated topic,
but also wondrous when implemented. For details, read
\hyperref[svn.webdav]{WebDAV and Autoversioning}.

Note that there are a number of other small tweaks one can make to
\texttt{mod\_dav\_svn} that perhaps do not merit extensive coverage. For
those interested, however, we provide a complete list of all
\texttt{httpd.conf} directives to which \texttt{mod\_dav\_svn} responds
in \hyperref[svn.serverconfig.httpd.ref.mod_dav_svn]{mod\_dav\_svn
configuration directives}.

\subsection{Subversion Apache HTTP Server Configuration
Reference}\label{svn.serverconfig.httpd.ref}

In the previous sections, we\textquotesingle ve mentioned numerous
directives that administrators can use in their \texttt{httpd.conf}
files to enable and configure their Subversion server offering,
introducing each directive as we also introduce the functionality it
toggles. In this section, we\textquotesingle ll quickly summarize
\emph{all} the configuration directives supported by both of the Apache
HTTP Server modules which are provided as part of the standard
Subversion distribution.

\subsubsection{mod\_dav\_svn configuration
directives}\label{svn.serverconfig.httpd.ref.mod_dav_svn}

The following configuration directives are recognized and supported by
Subversion\textquotesingle s Apache HTTP Server module,
\texttt{mod\_dav\_svn}.

\begin{description}
\item[\texttt{DAV\ svn}]
Must be included in any \texttt{Directory} or \texttt{Location} block
for a Subversion repository. It tells \texttt{httpd} to use the
Subversion backend for \texttt{mod\_dav} to handle all requests.
\item[\texttt{SVNActivitiesDB\ directory-path}]
Specifies the location in the filesystem where the activities database
should be stored. By default, \texttt{mod\_dav\_svn} creates and uses a
directory in the repository called \texttt{dav/activities.d}. The path
specified with this option must be an absolute path.

If specified for an \texttt{SVNParentPath} area, \texttt{mod\_dav\_svn}
appends the basename of the repository to the path specified here. For
example:

\begin{verbatim}
<Location /svn>
  DAV svn

  # any "/svn/foo" URL will map to a repository in 
  # /net/svn.nfs/repositories/foo
  SVNParentPath         "/net/svn.nfs/repositories"

  # any "/svn/foo" URL will map to an activities db in
  #  /var/db/svn/activities/foo
  SVNActivitiesDB       "/var/db/svn/activities"
</Location>
\end{verbatim}
\item[\texttt{SVNAdvertiseV2Protocol\ On\textbar{}Off}]
New to Subversion 1.7, this toggles whether \texttt{mod\_dav\_svn}
advertises its support for the new version of its HTTP protocol also
introduced in that version. Most admins will not wish to use this
directive (which is \texttt{On} by default), choosing instead to enjoy
the performance benefits that the new protocol offers. However, when
configuring a server as a write-through proxy to another server which
does not support the new protocol, set this directive\textquotesingle s
value to \texttt{Off}.
\item[\texttt{SVNAllowBulkUpdates\ On\textbar{}Off\textbar{}Prefer}]
Toggles support for all-inclusive responses to update-style requests.
Subversion clients use \texttt{REPORT} requests to get information about
directory tree checkouts and updates from \texttt{mod\_dav\_svn}. They
can ask the server to send that information in one of two ways: with the
entirety of the tree\textquotesingle s information in one massive (bulk)
response, or with a skelta (a skeletal representation of a tree delta)
which contains just enough information for the client to know what
\emph{additional} data to fetch from the server using subsequent
requests. When this directive is included with a value of \texttt{Off},
\texttt{mod\_dav\_svn} will only ever respond to these \texttt{REPORT}
requests with skelta responses, regardless of the type of responses
requested by the client.

The default value of this directive is \texttt{On}, which permits the
server to reply to update requests using the style of response (bulk or
skelta) requested by the client. Beginning in Subversion 1.8, this
directive also accepts a value of \texttt{Prefer}, which is similar to
\texttt{On} but additionally causes the server to announce to clients
that it \emph{prefers} to handle bulk update requests.

Most folks won\textquotesingle t need to use this directive at all. It
primarily exists for administrators who wish---for security or auditing
reasons---to force Subversion clients to fetch individually all the
files and directories needed for updates and checkouts, thus leaving an
audit trail of \texttt{GET} and \texttt{PROPFIND} requests in
Apache\textquotesingle s logs.
\item[\texttt{SVNAutoversioning\ On\textbar{}Off}]
When its value is \texttt{On}, allows write requests from WebDAV clients
to result in automatic commits. A generic log message is auto-generated
and attached to each revision. If you enable autoversioning,
you\textquotesingle ll likely want to set
\texttt{ModMimeUsePathInfo\ On} so that \texttt{mod\_mime} can set
\texttt{svn:mime-type} to the correct MIME type automatically (as best
as \texttt{mod\_mime} is able to, of course). For more information, see
\hyperref[svn.webdav]{WebDAV and Autoversioning}. The default value of
this directive is \texttt{Off}.
\item[\texttt{SVNCacheFullTexts\ On\textbar{}Off}]
When set to \texttt{On}, this tells Subversion to cache content
fulltexts if sufficient in-memory cache is available, which could offer
a significant performance benefit to the server. (See also the
\texttt{SVNInMemoryCacheSize} directive.) The default value of this
directive is \texttt{Off}.
\item[\texttt{SVNCacheTextDeltas\ On\textbar{}Off}]
When set to \texttt{On}, this tells Subversion to cache content deltas
if sufficient in-memory cache is available, which could offer a
significant performance benefit to the server. (See also the
\texttt{SVNInMemoryCacheSize} directive.) The default value of this
directive is \texttt{Off}.
\item[\texttt{SVNCompressionLevel\ level}]
Specifies the compression level used when sending file content over the
network. A value of \texttt{0} disables compression altogether, and
\texttt{9} is the maximum value. \texttt{5} is the default value.
\item[\texttt{SVNHooksEnv\ file-path}]
Specifies the location of the Subversion repository hook script
environment configuration file. This file is used to describe the
initial environment in which repository hook scripts are executed. For
more on this feature, see
\hyperref[svn.reposadmin.hooks.configuration]{Hook script environment
configuration}.
\item[\texttt{SVNIndexXSLT\ directory-path}]
Specifies the URI of an XSL transformation for directory indexes. This
directive is optional.
\item[\texttt{SVNInMemoryCacheSize\ size}]
Specifies the maximum size (in kbytes) per process of
Subversion\textquotesingle s in-memory object cache. The default value
is \texttt{16384}; use a value of \texttt{0} to deactivate this cache
altogether.
\item[\texttt{SVNListParentPath\ On\textbar{}Off}]
When set to \texttt{On}, allows a \texttt{GET} of
\texttt{SVNParentPath}, which results in a listing of all repositories
under that path. The default setting is \texttt{Off}.
\item[\texttt{SVNMasterURI\ url}]
Specifies a URI to the master Subversion repository (used for a
write-through proxy).
\item[\texttt{SVNMasterVersion\ X.Y}]
Specifies the release version number of the Subversion instance which is
serving the master repository (used for a write-through proxy).
\item[\texttt{SVNParentPath\ directory-path}]
Specifies the location in the filesystem of a parent directory whose
child directories are Subversion repositories. In a configuration block
for a Subversion repository, either this directive or \texttt{SVNPath}
must be present, but not both.
\item[\texttt{SVNPath\ directory-path}]
Specifies the location in the filesystem for a Subversion
repository\textquotesingle s files. In a configuration block for a
Subversion repository, either this directive or \texttt{SVNParentPath}
must be present, but not both.
\item[\texttt{SVNPathAuthz\ On\textbar{}Off\textbar{}short\_circuit}]
Controls path-based authorization by enabling subrequests (\texttt{On}),
disabling subrequests (\texttt{Off}; see
\hyperref[svn.serverconfig.httpd.authz.pathauthzoff]{Disabling
path-based checks}), or querying \texttt{mod\_authz\_svn} directly
(\texttt{short\_circuit}). The default value of this directive is
\texttt{On}.
\item[\texttt{SVNReposName\ name}]
Specifies the name of a Subversion repository for use in
\texttt{HTTP\ GET} responses. This value will be prepended to the title
of all directory listings (which are served when you navigate to a
Subversion repository with a web browser). This directive is optional.

Subversion will not use the repository name as configured via this
directive when trying to match rules in access control files. The
repository names used in that file\textquotesingle s syntax are always
derived from the repository URL. See
\hyperref[svn.serverconfig.pathbasedauthz.getting-started]{Getting
Started with Path-Based Access Control} for details.
\item[\texttt{SVNSpecialURI\ component}]
Specifies the URI component (namespace) for special Subversion
resources. The default is \texttt{!svn}, and most administrators will
never use this directive. Set this only if there is a pressing need to
have a file named \texttt{!svn} in your repository. If you change this
on a server already in use, it will break all of the outstanding working
copies, and your users will hunt you down with pitchforks and flaming
torches.
\item[\texttt{SVNUseUTF8\ On\textbar{}Off}]
When set to \texttt{On}, \texttt{mod\_dav\_svn} will communicate with
hook scripts using repository root paths encoded in UTF-8, and will
expect those scripts to likewise generate output (such as error
messages) encoded in UTF-8. The default value of this option is
\texttt{Off}, which means that \texttt{mod\_dav\_svn} assumes a 7-bit
ASCII encoding for its hook script interactions. This option is
available as of Subversion 1.8.

Administrators should ensure that the character set and encoding
expectations of hook scripts match all the ways they might be invoked.
For example, if one repository is served by both \texttt{httpd} and
\texttt{svnserve}, \texttt{svnserve} should also be configured to use
UTF-8 (by setting an appropriate locale in its environment) if this
option is enabled for \texttt{mod\_dav\_svn}. Also, local filesystem
paths containing non-ASCII characters which will be accessed by those
scripts (such as repository root paths) must be properly encoded in the
filesystem to match the scripts\textquotesingle{} expectations.
\end{description}

\subsubsection{mod\_authz\_svn configuration
directives}\label{svn.serverconfig.httpd.ref.mod_authz_svn}

The following configuration directives are provided by
\texttt{mod\_authz\_svn}, Subversion\textquotesingle s path-based
authorization Apache HTTP Server module. For an in-depth description of
using path-based authorization in Subversion, see
\hyperref[svn.serverconfig.pathbasedauthz]{Path-Based Authorization}.

\begin{description}
\item[\texttt{AuthzForceUsernameCase\ Upper\textbar{}Lower}]
Set to \texttt{Upper} or \texttt{Lower} to perform case conversion of
the specified sort on the authenticated username before checking it for
authorization. While usernames are compared in a case-sensitive fashion
against those referenced in the authorization rules file, this directive
can at least normalize variably-cased usernames into something
consistent.
\item[\texttt{AuthzSVNAccessFile\ file-path}]
Consult \textless file-path\textgreater{} for access rules describing
the permissions for paths in Subversion repository. In a configuration
block for a Subversion repository or a colletion of repositories, either
this directive or \texttt{AuthzSVNReposRelativeAccessFile} can be
present, but not both.

Beginning with Subversion 1.8, \texttt{AuthzSVNAccessFile} accepts a URL
to a file stored inside a Subversion repository---either the same
repository to which the access rules are being applied, or an entirely
different repository.
\item[\texttt{AuthzSVNAnonymous\ On\textbar{}Off}]
Set to \texttt{Off} to disable two special-case behaviours of this
module: interaction with the \texttt{Satisfy\ Any} directive and
enforcement of the authorization policy even when no \texttt{Require}
directives are present. The default value of this directive is
\texttt{On}.
\item[\texttt{AuthzSVNAuthoritative\ On\textbar{}Off}]
Set to \texttt{Off} to allow access control to be passed along to lower
modules. The default value of this directive is \texttt{On}.
\item[\texttt{AuthzSVNNoAuthWhenAnonymousAllowed\ On\textbar{}Off}]
Set to \texttt{On} to suppress authentication and authorization for
requests which anonymous users are allowed to perform. The default value
of this directive is \texttt{On}.
\item[\texttt{AuthzSVNReposRelativeAccessFile\ file-path}]
Consult \textless file-path\textgreater{} for access rules describing
the permissions for paths in Subversion repository. Unlike
\texttt{AuthzSVNAccessFile}, the path specified for
\texttt{AuthzSVNReposRelativeAccessFile} is relative to the
\texttt{conf/} directory in the repository on filesystem. In other
words, the \textless file-path\textgreater{} specifies a per repository
file that must be accessible by the relative path for all repositories
in a configuration block. In a configuration block for a Subversion
repository or a collection of repositories either this directive or
\texttt{AuthzSVNAccessFile} can be present, but not both. This option is
available as of Subversion 1.7.

Beginning with Subversion 1.8, \texttt{AuthzSVNReposRelativeAccessFile}
accepts a URL to a file stored inside a Subversion repository---either
the same repository to which the access rules are being applied, or an
entirely different repository.
\end{description}

\section{Path-Based
Authorization}\label{svn.serverconfig.pathbasedauthz}

Both Apache and \texttt{svnserve} are capable of granting (or denying)
permissions to users. Typically this is done over the entire repository:
a user can read the repository (or not), and she can write to the
repository (or not).

It\textquotesingle s also possible, however, to define finer-grained
access rules. One set of users may have permission to write to a certain
directory in the repository, but not others; another directory might not
even be readable by all but a few special people. It\textquotesingle s
even possible to restrict access on a per file basis.

Both Subversion servers use a common file format to describe these
path-based access rules. In this section, we will explain that file
format, as well how to configure your Subversion server to use it for
managing path-based authorization.

Do You Really Need Path-Based Access Control?

A lot of administrators setting up Subversion for the first time tend to
jump into path-based access control without giving it a lot of thought.
The administrator usually knows which teams of people are working on
which projects, so it\textquotesingle s easy to jump in and grant
certain teams access to certain directories and not others. It seems
like a natural thing, and it appeases the
administrator\textquotesingle s desire to maintain tight control of the
repository.

Note, though, that there are often invisible (and visible!) costs
associated with this feature. In the visible category, the server needs
to do a lot more work to ensure that the user has the right to read or
write each specific path; in certain situations, there\textquotesingle s
very noticeable performance loss. In the invisible category, consider
the culture you\textquotesingle re creating. Most of the time, while
certain users \emph{shouldn\textquotesingle t} be committing changes to
certain parts of the repository, that social contract
doesn\textquotesingle t need to be technologically enforced. Teams can
sometimes spontaneously collaborate with each other; someone may want to
help someone else out by committing to an area she
doesn\textquotesingle t normally work on. By preventing this sort of
thing at the server level, you\textquotesingle re setting up barriers to
unexpected collaboration. You\textquotesingle re also creating a bunch
of rules that need to be maintained as projects develop, new users are
added, and so on. It\textquotesingle s a bunch of extra work to
maintain.

Remember that this is a version control system! Even if somebody
accidentally commits a change to something she
shouldn\textquotesingle t, it\textquotesingle s easy to undo the change.
And if a user commits to the wrong place with deliberate malice,
it\textquotesingle s a social problem anyway, and that the problem needs
to be dealt with outside Subversion.

So, before you begin restricting users\textquotesingle{} access rights,
ask yourself whether there\textquotesingle s a real, honest need for
this, or whether it\textquotesingle s just something that ``sounds
good'' to an administrator. Decide whether it\textquotesingle s worth
sacrificing some server speed, and remember that there\textquotesingle s
very little risk involved; it\textquotesingle s bad to become dependent
on technology as a crutch for social problems.\footnote{A common theme
  in this book!}

As an example to ponder, consider that the Subversion project itself has
always had a notion of who is allowed to commit where, but
it\textquotesingle s always been enforced socially. This is a good model
of community trust, especially for open source projects. Of course,
sometimes there \emph{are} truly legitimate needs for path-based access
control; within corporations, for example, certain types of data really
can be sensitive, and access needs to be genuinely restricted to small
groups of people.

\subsection{Getting Started with Path-Based Access
Control}\label{svn.serverconfig.pathbasedauthz.getting-started}

Subversion offers path-based access control in Apache via the
\texttt{mod\_authz\_svn} module, which must be loaded using the
\texttt{LoadModule} directive in \texttt{httpd.conf} in the same fashion
that \texttt{mod\_dav\_svn} itself is loaded. To enable the use of this
module for your repositories, you\textquotesingle ll add the
\texttt{AuthzSVNAccessFile} or \texttt{AuthzSVNReposRelativeAccessFile}
directives (again within the \texttt{httpd.conf} file) pointing to your
own access rules file. (For a full explanation, see
\hyperref[svn.serverconfig.httpd.authz.perdir]{Per-directory access
control}.)

To configure path-based authorization in \texttt{svnserve}, simply point
the \texttt{authz-db} configuration variable (within your
\texttt{svnserve.conf} file) to your access rules file.

Once your server knows where to look for your access rules,
it\textquotesingle s time to define those rules.

The syntax of the Subversion access file is the same familiar one used
by \texttt{svnserve.conf} and the runtime configuration files. Lines
that start with a hash (\texttt{\#}) are ignored. In its simplest form,
each section names a versioned path and, optionally, the repository in
which that path is found. In other words, except for a few reserved
sections, section names are of one of two forms: either
\texttt{{[}repos-name:path{]}} or \texttt{{[}path{]}} when
\texttt{AuthzSVNAccessFile} is used. If you configured per repository
access files via \texttt{AuthzSVNReposRelativeAccessFile} directive, you
should always use \texttt{{[}path{]}} form only. (Subversion 1.10 added a
third, wildcard-aware form of section name; see
\hyperref[svn.serverconfig.pathbasedauthz.wildcards]{Path Patterns with
Wildcards}.) Authenticated usernames
are the option names within each section, and an
option\textquotesingle s value describes that user\textquotesingle s
level of access to the repository path: either \texttt{r} (read-only) or
\texttt{rw} (read/write). If the user is not mentioned at all, no access
is allowed.

Paths used in access file sections must be specified using
Subversion\textquotesingle s ``internal style'', which mostly just means
that they are encoded in UTF-8 and use forward slash (\texttt{/})
characters as directory separators (even on Windows systems). Note also
that these paths do not employ any character escaping mechanism (such as
URI-encoding)---spaces in path names should be represented exactly as
such in access file section names
(\texttt{{[}repos-name:path\ with\ spaces{]}}, e.g.)

Here\textquotesingle s a simple example demonstrating a piece of the
access configuration which grants read access Sally, and read/write
access to Harry, for the path \texttt{/branches/calc/bug-142} (and all
its children) in the repository \texttt{calc}:

\begin{verbatim}
[calc:/branches/calc/bug-142]
harry = rw
sally = r
\end{verbatim}

Prior to version 1.7, Subversion treated repository names and paths in a
case-insensitive fashion for the purposes of access control, converting
them to lower case internally before comparing them against the contents
of your access file. It now does these comparisons case-sensitively. If
you upgraded to Subversion 1.7 from an older version, you should review
your access files for case correctness.

The name of a repository as evaluated by the authorization subsystem is
derived directly from the repository\textquotesingle s path. Exactly how
this happens differs between the two server options.
\texttt{mod\_dav\_svn} uses only the basename of the
repository\textquotesingle s root URL\footnote{Any human-readable name
  for a repository configured via the \texttt{SVNReposName}
  \texttt{httpd.conf} directive will be ignored by the authorization
  subsystem. Your access control file sections must refer to
  repositories by their server-sensitive paths as previously described.},
while \texttt{svnserve} uses the entire relative path from the serving
root (as determined by its \texttt{-\/-root} (\texttt{-r}) command-line
option) to the repository.

The differences in the ways that a repository\textquotesingle s name is
determined by each of \texttt{mod\_dav\_svn} and \texttt{svnserve} can
cause problems when trying to serve a repository via both servers
simultaneously. Naturally, an administrator would prefer to point both
servers\textquotesingle{} configurations toward a common access file.
However, for this to work, you must ensure that the repository name
portion of the file\textquotesingle s section names are compatible with
each server\textquotesingle s idea of what the repository name should
be---for example, by configuring \texttt{svnserve}\textquotesingle s
root to be the same as \texttt{mod\_dav\_svn}\textquotesingle s
configured \texttt{SVNParentPath}, or by using a different access file
per repository so that section names needn\textquotesingle t reference
the repository at all.

If you\textquotesingle re using the \texttt{SVNParentPath} directive,
it\textquotesingle s important to specify the repository names in your
sections. If you omit them, a section such as \texttt{{[}/some/dir{]}}
will match the path \texttt{/some/dir} in \emph{every} repository. If
you\textquotesingle re using the \texttt{SVNPath} directive, however,
it\textquotesingle s fine to provide only paths in your sections---after
all, there\textquotesingle s only one repository.

Permissions are inherited from a path\textquotesingle s parent
directory. That means we can specify a subdirectory with a different
access policy for Sally. Let\textquotesingle s continue our previous
example, and grant Sally write access to a child of the branch that
she\textquotesingle s otherwise permitted only to read:

\begin{verbatim}
[calc:/branches/calc/bug-142]
harry = rw
sally = r

# give sally write access only to the 'testing' subdir
[calc:/branches/calc/bug-142/testing]
sally = rw
\end{verbatim}

Now Sally can write to the \texttt{testing} subdirectory of the branch,
but can still only read other parts. Harry, meanwhile, continues to have
complete read/write access to the whole branch.

It\textquotesingle s also possible to explicitly deny permission to
someone via inheritance rules, by setting the username variable to
nothing:

\begin{verbatim}
[calc:/branches/calc/bug-142]
harry = rw
sally = r

[calc:/branches/calc/bug-142/secret]
harry =
\end{verbatim}

In this example, Harry has read/write access to the entire
\texttt{bug-142} tree, but has absolutely no access at all to the
\texttt{secret} subdirectory within it.

The thing to remember is that the most specific path always matches
first. The server tries to match the path itself, and then the parent of
the path, then the parent of that, and so on. The net effect is that
mentioning a specific path in the access file will always override any
permissions inherited from parent directories.

Similarly, sections that specify a repository name have precedence over
those that don\textquotesingle t: if both \texttt{{[}calc:/some/path{]}}
and \texttt{{[}/some/path{]}} are present, the former will be used and
the latter ignored for \texttt{calc}.

By default, nobody has any access to any repository at all. That means
that if you\textquotesingle re starting with an empty file,
you\textquotesingle ll probably want to give at least read permission to
all users at the roots of the repositories. You can do this by using the
asterisk variable (\texttt{*}), which means ``all users'':

\begin{verbatim}
[/]
* = r
\end{verbatim}

This is a common setup; notice that no repository name is mentioned in
the section name. This makes all repositories world-readable to all
users. Once all users have read access to the repositories, you can give
explicit \texttt{rw} permission to certain users on specific
subdirectories within specific repositories.

Note that while all of the previous examples use directories,
that\textquotesingle s only because defining access rules on directories
is the most common case. You may similarly restrict access on file
paths, too.

\begin{verbatim}
[calendar:/projects/calendar/manager.ics]
harry = rw
sally = r
\end{verbatim}

\subsection{Access Control
Groups}\label{svn.serverconfig.pathbasedauthz.groups}

The access file also allows you to define whole groups of users, much
like the Unix \texttt{/etc/group} file. To do this, create a
\texttt{groups} section in your access file, and then describe your
groups within that section: each variable\textquotesingle s name defines
the name of the group, and its value is a comma-delimited list of
usernames which are part of that group.

\begin{verbatim}
[groups]
calc-developers = harry, sally, joe
paint-developers = frank, sally, jane
everyone = harry, sally, joe, frank, jane
\end{verbatim}

Groups can be granted access control just like users. Distinguish them
with an ``at sign'' (\texttt{@}) prefix:

\begin{verbatim}
[calc:/projects/calc]
@calc-developers = rw

[paint:/projects/paint]
jane = r
@paint-developers = rw
\end{verbatim}

Another important fact is that group permissions are not overridden by
individual user permissions. Rather, the \emph{combination} of all
matching permissions is granted. In the prior example, Jane is a member
of the \texttt{paint-developers} group, which has read/write access.
Combined with the \texttt{jane\ =\ r} rule, this still gives Jane
read/write access. Permissions for group members can only be extended
beyond the permissions the group already has. Restricting users who are
part of a group to less than their group\textquotesingle s permissions
is impossible.

Groups can also be defined to contain other groups:

\begin{verbatim}
[groups]
calc-developers = harry, sally, joe
paint-developers = frank, sally, jane
everyone = @calc-developers, @paint-developers
\end{verbatim}

\subsection{Username
Aliases}\label{svn.serverconfig.pathbasedauthz.aliases}

Some authentication systems expect and carry relatively short usernames
of the sorts we\textquotesingle ve been describing
here---\texttt{harry}, \texttt{sally}, \texttt{joe}, and so on. But
other authentication systems---such as those which use LDAP stores or
SSL client certificates---may carry much more complex usernames. For
example, Harry\textquotesingle s username in an LDAP-protected system
might be \texttt{CN=Harold\ Hacker,OU=Engineers,DC=red-bean,DC=com}.
With usernames like that, the access file can become quite bloated with
long or obscure usernames that are easy to mistype.

Fortunately, Subversion 1.5 introduced username aliases to the access
file syntax. Username aliases allow you to have to type the correct
complex username only once, in a statement which assigns to it a more
easily digestable alias.

Username aliases are defined in the special \texttt{aliases} section of
the access file, with each variable name in that section defining an
alias, and the value of those variables carrying the real Subversion
username which is being aliased.

\begin{verbatim}
[aliases]
harry = CN=Harold Hacker,OU=Engineers,DC=red-bean,DC=com
sally = CN=Sally Swatterbug,OU=Engineers,DC=red-bean,DC=com
joe = CN=Gerald I. Joseph,OU=Engineers,DC=red-bean,DC=com
…
\end{verbatim}

Once you\textquotesingle ve defined a set of aliases, you can refer to
the users elsewhere in the access file via their aliases in all the same
places you could have instead used their actual usernames. Simply
prepend an ampersand to the alias to distinguish it from a regular
username:

\begin{verbatim}
[groups]
calc-developers = &harry, &sally, &joe
paint-developers = &frank, &sally, &jane
everyone = @calc-developers, @paint-developers
\end{verbatim}

You might also choose to use aliases if your users\textquotesingle{}
usernames change frequently. Doing so allows you to need to update only
the aliases table when these username changes occur, instead of doing
global search-and-replace operations on the whole access file.

\subsection{Advanced Access Control
Features}\label{svn.serverconfig.pathbasedauthz.authclass-tokens}

Beginning with Subversion 1.5, the access file syntax also supports some
``magic'' tokens for helping you to make rule assignments based on the
user\textquotesingle s authentication class. One such token is the
\texttt{\$authenticated} token. Use this token where you would otherwise
specify a username, alias, or group name in your authorization rules to
declare the permissions granted to any user who has authenticated with
any username at all. Similarly employed is the \texttt{\$anonymous}
token, except that it matches everyone who has \emph{not} authenticated
with a username.

\begin{verbatim}
[calendar:/projects/calendar]
$anonymous = r
$authenticated = rw
\end{verbatim}

Another handy bit of access file syntax magic is the use of the tilde
(\texttt{\textasciitilde{}}) character as an exclusion marker. In your
authorization rules, prefixing a username, alias, group name, or
authentication class token with a tilde character will cause Subversion
to apply the rule to users who do \emph{not} match the rule. Though
somewhat unnecessarily obfuscated, the following block is equivalent to
the one in the previous example:

\begin{verbatim}
[calendar:/projects/calendar]
~$authenticated = r
~$anonymous = rw
\end{verbatim}

A less obvious example might be as follows:

\begin{verbatim}
[groups]
calc-developers = &harry, &sally, &joe
calc-owners = &hewlett, &packard
calc = @calc-developers, @calc-owners

# Any calc participant has read-write access...
[calc:/projects/calc]
@calc = rw

# ...but only allow the owners to make and modify release tags.
[calc:/projects/calc/tags]
~@calc-owners = r
\end{verbatim}

\subsection{Path Patterns with
Wildcards}\label{svn.serverconfig.pathbasedauthz.wildcards}

The access rules described so far each name a single, literal path. As of
Subversion 1.10, a rule\textquotesingle s path may instead contain
wildcards, letting a single rule apply to many paths at once. A wildcard
rule is written in a section whose name carries a \texttt{:glob:} marker
immediately after the opening bracket---either
\texttt{{[}:glob:repos-name:/path{]}} or \texttt{{[}:glob:/path{]}}, the
latter (without a repository name) applying across all repositories.

Two wildcards are recognized within the path:

\begin{description}
\item[\texttt{*}]
Matches exactly one arbitrary path segment. It never matches a slash, so
it stays within a single level of the directory tree.
\item[\texttt{**}]
Matches an arbitrary number of path segments (zero or more) separated by
slashes. Written as \texttt{/**/} it can match zero intervening segments,
so \texttt{/*/**/*} matches any path of at least two segments.
\end{description}

Classic shell-style patterns work within a single segment as well---for
example, \texttt{*.c} matches all the C source files in a directory---and
a literal asterisk can be matched by escaping it with a backslash.

For instance, to grant a group read/write access to the \texttt{tags}
directory of \emph{every} project in the \texttt{calc} repository, no
matter how deeply nested it is:

\begin{verbatim}
[:glob:calc:/**/tags]
@calc-developers = rw
\end{verbatim}

Or to make every C source file throughout a repository read-only for a
group of reviewers:

\begin{verbatim}
[:glob:calc:/**/*.c]
@reviewers = r
\end{verbatim}

Literal (non-wildcard) rules continue to take precedence over wildcard
rules, so you can still carve out exceptions with an ordinary section. To
avoid ambiguity, Subversion refuses to load an access file in which a
wildcard rule and a literal rule resolve to the very same path; such
collisions are reported as an error rather than silently resolved.

\subsection{Some Gotchas with Access
Control}\label{svn.serverconfig.pathbasedauthz.gotchas}

If you\textquotesingle re using Apache as your Subversion server and
have made certain subdirectories of your repository unreadable to
certain users, you need to be aware of a possible nonoptimal behavior
with \texttt{svn\ checkout}.

Depending on which HTTP communication library the Subversion client is
using, it may request that the entire payload of a checkout or update be
delivered in a single (often large) response to the primary
checkout/update request. When this happens, this single request is the
\emph{only} opportunity Apache has to demand user authentication. This
has some odd side effects. For example, if a certain subdirectory of the
repository is readable only by user Sally, and user Harry checks out a
parent directory, his client will respond to the initial authentication
challenge as Harry. As the server generates the large response,
there\textquotesingle s no way it can resend an authentication challenge
when it reaches the special subdirectory; thus the subdirectory is
skipped altogether, rather than asking the user to reauthenticate as
Sally at the right moment.

In a similar way, if the root of the repository is anonymously
world-readable, the entire checkout will be done without
authentication---again, skipping the unreadable directory, rather than
asking for authentication partway through.\footnote{For more on this,
  see the blog post \emph{Authz and Anon Authn Agony} at
  \href{https://subversion.apache.org/blog/2007-03-27-authz-and-anon-authn-agony.html}{}.}

\section{High-level Logging}\label{svn.serverconfig.operational-logging}

Both the Apache \texttt{httpd} and \texttt{svnserve} Subversion servers
provide support for high-level logging of Subversion operations.
Configuring each of the server options to provide this level of logging
is done differently, of course, but the output from each is designed to
conform to a uniform syntax.

To enable high-level logging in \texttt{svnserve}, you need only use the
\texttt{-\/-log-file} command-line option when starting the server,
passing as the value to the option the file to which \texttt{svnserve}
should write its log output.

\begin{verbatim}
$ svnserve -d -r /path/to/repositories --log-file /var/log/svn.log
\end{verbatim}

Enabling the same in Apache is a bit more involved, but is essentially
an extension of Apache\textquotesingle s stock log output configuration
mechanisms---see \hyperref[svn.serverconfig.httpd.extra.logging]{Apache
logging} for details.

The following is a list of Subversion action log messages produced by
its high-level logging mechanism, followed by one or more examples of
the log message as it appears in the log output.

\begin{description}
\item[Checkout or export] ~ 
\begin{verbatim}
checkout-or-export /path r62 depth=infinity
\end{verbatim}
\item[Commit] ~ 
\begin{verbatim}
commit harry r100
\end{verbatim}
\item[Diffs] ~ 
\begin{verbatim}
diff /path r15:20 depth=infinity ignore-ancestry
diff /path1@15 /path2@20 depth=infinity ignore-ancestry
\end{verbatim}
\item[Fetch a directory] ~ 
\begin{verbatim}
get-dir /trunk r17 text
\end{verbatim}
\item[Fetch a file] ~ 
\begin{verbatim}
get-file /path r20 props
\end{verbatim}
\item[Fetch a file revision] ~ 
\begin{verbatim}
get-file-revs /path r12:15 include-merged-revisions
\end{verbatim}
\item[Fetch merge information] ~ 
\begin{verbatim}
get-mergeinfo (/path1 /path2)
\end{verbatim}
\item[Lock] ~ 
\begin{verbatim}
lock /path steal
\end{verbatim}
\item[Log] ~ 
\begin{verbatim}
log (/path1,/path2,/path3) r20:90 discover-changed-paths revprops=()
\end{verbatim}
\item[Replay revisions (svnsync)] ~ 
\begin{verbatim}
replay /path r19
\end{verbatim}
\item[Revision property change] ~ 
\begin{verbatim}
change-rev-prop r50 propertyname
\end{verbatim}
\item[Revision property list] ~ 
\begin{verbatim}
rev-proplist r34
\end{verbatim}
\item[Status] ~ 
\begin{verbatim}
status /path r62 depth=infinity
\end{verbatim}
\item[Switch] ~ 
\begin{verbatim}
switch /pathA /pathB@50 depth=infinity
\end{verbatim}
\item[Unlock] ~ 
\begin{verbatim}
unlock /path break
\end{verbatim}
\item[Update] ~ 
\begin{verbatim}
update /path r17 send-copyfrom-args
\end{verbatim}
\end{description}

As a convenience to administrators who wish to post-process their
Subversion high-level logging output (perhaps for reporting or analysis
purposes), Subversion source code distributions provide a Python module
(located at \texttt{tools/server-side/svn\_server\_log\_parse.py}) which
can be used to parse Subversion\textquotesingle s log output.

\section{Server Optimization}\label{svn.serverconfig.optimization}

Part of the due diligence when offering a service such as a Subversion
server involves capacity planning and performance tuning. Subversion
doesn\textquotesingle t tend to be particularly greedy in terms of
server resources such as CPU cycles and memory, but any service can
benefit from optimizations, especially when usage of the service
skyrockets\footnote{In Subversion\textquotesingle s case, the
  skyrocketing affect is, of course, due to its cool name. Well, that
  and its popularity, reliability, ease of use\ldots.}. In this section,
we\textquotesingle ll discuss some ways you can tweak your Subversion
server configuration to offer even better performance and scalability.

\subsection{Data Caching}\label{svn.serverconfig.optimization.caching}

Generally speaking, the most expensive part of a Subversion
server\textquotesingle s job is fetching data from the repository.
Subversion 1.6 attempted to offset this cost by introducing some
in-memory caching of certain classes of data read from the repository.
But Subversion 1.7 takes this a step further, not only caching the
results of some of the more costly operations, but also by providing in
each of the available servers the means by which fine-tune the size and
some behaviors of the cache.

For \texttt{svnserve}, you can specify the size of the cache using the
\texttt{-\/-memory-cache-size} (\texttt{-M}) command-line option. You
can also dictate whether \texttt{svnserve} should attempt to cache
content fulltexts and deltas via the boolean
\texttt{-\/-cache-fulltexts} and \texttt{-\/-cache-txdeltas} options,
respectively.

\begin{verbatim}
$ svnserve -d -r /path/to/repositories \
           --memory-cache-size 1024 \
           --cache-txdeltas yes \
           --cache-fulltexts yes
…
$
\end{verbatim}

\texttt{mod\_dav\_svn} provides the same degree of cache configurability
via \texttt{httpd.conf} directives. The \texttt{SVNInMemoryCacheSize},
\texttt{SVNCacheFullTexts}, and \texttt{SVNCacheTextDeltas} directives
may be used at the server configuration level to control
Subversion\textquotesingle s data cache characteristics:

\begin{verbatim}
<IfModule dav_svn_module>
  # Enable a 1 Gb Subversion data cache for both fulltext and deltas.
  SVNInMemoryCacheSize 1048576
  SVNCacheTextDeltas On
  SVNCacheFullTexts On
</IfModule>
\end{verbatim}

So what settings should you use? Certainly you need to consider what
resources are available on your server. To get any benefit out of the
cache at all, you\textquotesingle ll probably want to let the cache be
at least large enough to hold all the files which are most commonly
accessed in your repository (for example, your project\textquotesingle s
\texttt{trunk} directory tree).

Setting the memory cache size to \texttt{0} will disable this enhanced
caching mechanism and cause Subversion to fall back to using the older
cache mechanisms introduced in Subversion 1.6.

Currently, only repositories which make use of the FSFS backend data
store make use of this data caching functionality.

\subsection{Network Compression of
Data}\label{svn.serverconfig.optimization.compression}

Compressing the data transmitted across the wire can greatly reduce the
size of those network transmissions, but comes at the cost of server
(and client) CPU cycles. Depending on your server\textquotesingle s CPU
capacity, the typical access patterns of the clients who use your
servers, and the bandwidth of the networks between them, you might wish
to fine tune just how hard your server will work to compress the data it
sends across the wire. To assist with this fine tuning process,
Subversion 1.7 offers the \texttt{-\/-compression} (\texttt{-c}) option
to \texttt{svnserve} and the \texttt{SVNCompressionLevel} directive for
\texttt{mod\_dav\_svn}. Both accept a value which is an integer between
0 and 9 (inclusive), where 9 offers the best compression of wire data,
and 0 disables compression altogether.

For example, on a local area network (LAN) with 1-Gigabit connections,
it might not make sense to have the server compress its network
transmissions (which also forces the clients to decompress them), as the
network itself is so fast that users won\textquotesingle t really
benefit from the smaller overall network payload. On the other hand,
servers which are accessed primarily by clients with low-bandwidth
connections would be doing those clients a favor by minimizing the
overall size of its network communications.

\section{Supporting Multiple Repository Access
Methods}\label{svn.serverconfig.multimethod}

You\textquotesingle ve seen how a repository can be accessed in many
different ways. But is it possible---or safe---for your repository to be
accessed by multiple methods simultaneously? The answer is yes, provided
you use a bit of foresight.

At any given time, these processes may require read and write access to
your repository:

\begin{itemize}
\item
  Regular system users using a Subversion client (as themselves) to
  access the repository directly via \texttt{file://} URLs
\item
  Regular system users connecting to SSH-spawned private
  \texttt{svnserve} processes (running as themselves), which access the
  repository
\item
  An \texttt{svnserve} process---either a daemon or one launched by
  \texttt{inetd}---running as a particular fixed user
\item
  An Apache \texttt{httpd} process, running as a particular fixed user
\end{itemize}

The most common problem administrators run into is repository ownership
and permissions. Does every process (or user) in the preceding list have
the rights to read and write the repository\textquotesingle s underlying
data files? Assuming you have a Unix-like operating system, a
straightforward approach might be to place every potential repository
user into a new \texttt{svn} group, and make the repository wholly owned
by that group. But even that\textquotesingle s not enough, because a
process may write to the database files using an unfriendly umask---one
that prevents access by other users.

So the next step beyond setting up a common group for repository users
is to force every repository-accessing process to use a sane umask. For
users accessing the repository directly, you can make the \texttt{svn}
program into a wrapper script that first runs \texttt{umask\ 002} and
then runs the real \texttt{svn} client program. You can write a similar
wrapper script for the \texttt{svnserve} program, and add a
\texttt{umask\ 002} command to Apache\textquotesingle s own startup
script, \texttt{apachectl}. For example:

\begin{verbatim}
$ cat /usr/bin/svn

#!/bin/sh

umask 002
/usr/bin/svn-real "$@"
\end{verbatim}

Another common problem is often encountered on Unix-like systems. As
your repository accumulates revisions, Subversion creates new files
within it. Even if the repository is wholly owned by the \texttt{svn}
group, these newly created files won\textquotesingle t necessarily be
owned by that same group, which then creates more permissions problems
for your users. A good workaround is to set the group SUID bit on the
repository\textquotesingle s \texttt{db} directory. This causes all
newly created files to have the same group owner as the parent
directory.

Once you\textquotesingle ve jumped through these hoops, your repository
should be accessible by all the necessary processes. It may seem a bit
messy and complicated, but the problems of having multiple users sharing
write access to common files are classic ones that are not often
elegantly solved.

Fortunately, most repository administrators will never \emph{need} to
have such a complex configuration. Users who wish to access repositories
that live on the same machine are not limited to using \texttt{file://}
access URLs---they can typically contact the Apache HTTP server or
\texttt{svnserve} using \texttt{localhost} for the server name in their
\texttt{http://} or \texttt{svn://} URL. And maintaining multiple server
processes for your Subversion repositories is likely to be more of a
headache than necessary. We recommend that you choose a single server
that best meets your needs and stick with it!

The svn+ssh:// Server Checklist

It can be quite tricky to get a bunch of users with existing SSH
accounts to share a repository without permissions problems. If
you\textquotesingle re confused about all the things that you (as an
administrator) need to do on a Unix-like system, here\textquotesingle s
a quick checklist that resummarizes some of the topics discussed in this
section:

\begin{itemize}
\item
  All of your SSH users need to be able to read and write to the
  repository, so put all the SSH users into a single group.
\item
  Make the repository wholly owned by that group.
\item
  Set the group permissions to read/write.
\item
  Your users need to use a sane umask when accessing the repository, so
  make sure \texttt{svnserve} (\texttt{/usr/bin/svnserve}, or wherever
  it lives in \texttt{\$PATH}) is actually a wrapper script that runs
  \texttt{umask\ 002} and executes the real \texttt{svnserve} binary.
\item
  Take similar measures when using \texttt{svnlook} and
  \texttt{svnadmin}. Either run them with a sane umask or wrap them as
  just described.
\end{itemize}

